\chapimg{fig/sh10.jpg}
\IfFileExists{extra/lead10.tex}{\chapterlead{\subhd{After this chapter you can} define phenotyping traits for water use; compute water productivity, harvest index and a crop water stress index; rank genotypes under managed stress; and plan the volumetric scan and the plot design of a trial.}
}{}
\chapter{Phenotyping: the Deterministic Plant and the Volumetric Scan}\label{chap:10}
\markboth{Phenotyping: the Deterministic Plant and the Volumetric Scan}{Phenotyping: the Deterministic Plant and the Volumetric Scan}
\IfFileExists{extra/bg10.tex}{\section{Introduction: why phenotype water use}
Phenotyping is the measurement of the traits of a plant in its environment. For water use the useful traits are the yield, the water used, the way the canopy takes and loses water, and the way it responds to stress. The application treats phenotyping as a layer on top of the balance: the balance supplies the water, the trait tables supply the response of each genotype, and the comparison of genotypes is made on the same water.

\section{Traits and indices}
The main traits are the \emph{water productivity} $WP=Y/ET_a$, the \emph{harvest index} $HI=Y/B$ (grain over biomass), the \emph{transpiration efficiency} $TE=B/T$, the \emph{stomatal conductance} $g_s$ and the \emph{canopy temperature} $T_c$. The last gives the \kw{crop water stress index}

\begin{equation}CWSI=\frac{(T_c-T_a)-(T_c-T_a)_{LL}}{(T_c-T_a)_{UL}-(T_c-T_a)_{LL}},\quad 0\le CWSI\le1,\end{equation}


where LL and UL are the non-stressed and fully stressed baselines. A well-watered canopy runs 2 to 4 C below the air; a stressed one runs above it.

\section{Experimental design}
A trial is a set of plots in which genotype is the factor and the other conditions are the same. The \emph{randomised complete block} design gives every genotype one plot in every block, so that gradients of soil and slope are removed from the comparison. With $g$ genotypes and $r$ blocks the degrees of freedom of the error are $(g-1)(r-1)$, and the standard error of a genotype mean is $s/\sqrt r$. A trial of 20 genotypes in 4 blocks has 76 error degrees of freedom and 80 plots.

The least significant difference at $p=0.05$ is $LSD=t_{0.05,\nu}\sqrt{2s^2/r}$; with $s=0.6$ Mg\,ha$^{-1}$, $r=4$, $t\approx2.0$: $LSD=0.85$ Mg\,ha$^{-1}$.

\section{Volumetric scan and image traits}
The application's scan reads a series of images or point-cloud slices of a plant and returns the volume of the canopy and of the root system. The biomass is estimated by $B=\rho_{v}V$ where $\rho_v$ is the volumetric dry density (kg\,m$^{-3}$) of the tissue. A canopy volume of 0.18 m$^3$ per plant at $\rho_v=18$ kg\,m$^{-3}$ gives $B=3.24$ kg per plant; at 5 plants\,m$^{-2}$ the biomass is 162.0 Mg\,ha$^{-1}$.

}{}
\IfFileExists{extra/deep10.tex}{\section{The two phenotyping layers}
The application has two layers of phenotyping. The first, \emph{deterministic phenotyping}, is an exactly reproducible virtual experiment: a plant is grown by a model driven by the soil, the tillage, the weather, the irrigation, the mulch and the crop of the project, with a well-watered treatment and a drought treatment side by side. Nothing in it is random between runs, so every control of the project moves the numbers and two runs with the same inputs give the same plant. The second, \emph{volumetric phenotyping}, places a virtual scanner around that plant, in the way a stereo-camera rig looks at a seedling on a turntable, and measures the volume, the crown width, the height, the leaf area and the number of leaves from the point cloud that the scanner records. The difference between the two is the difference between the truth of the model and what an instrument would report.

\section{The scanner: views, registration and the iterative closest point}
Each view is a cloud of points in its own frame. The method follows Wang et al.\ (2025): structure from motion and multi-view stereo with six views at intervals of $60^\circ$ and calibration spheres placed beside the plant, used to join the views (self-registration), and a refinement by the iterative closest point (ICP) algorithm. The joining of two clouds is the search for the rotation $\mathbf R$ and translation $\mathbf t$ that minimise

\begin{equation}\min_{\mathbf R,\mathbf t}\sum_{i=1}^{n}\left\|\mathbf R\,\mathbf p_i+\mathbf t-\mathbf q_i\right\|^{2},\qquad \mathbf R=\mathbf V\,\mathrm{diag}\bigl(1,1,\det(\mathbf V\mathbf U^{\!\top})\bigr)\,\mathbf U^{\!\top},\quad \mathbf t=\bar{\mathbf q}-\mathbf R\bar{\mathbf p},\label{eq:ph-icp}\end{equation}
\par\noindent where $\mathbf p_i$ and $\mathbf q_i$ are corresponding points of the moving and the fixed cloud, $n$ the number of pairs, $\mathbf R$ the rotation matrix, $\mathbf t$ the translation vector, $\bar{\mathbf p}$ and $\bar{\mathbf q}$ the centroids, and $\mathbf U\boldsymbol\Sigma\mathbf V^{\!\top}$ the singular-value decomposition of the covariance $\sum(\mathbf p_i-\bar{\mathbf p})(\mathbf q_i-\bar{\mathbf q})^{\!\top}$ (the determinant factor excludes reflections).



The calibration spheres give the first correspondences (their centres are seen in every view); ICP then pairs every point with its nearest neighbour in the other cloud, solves Eq.~(\ref{eq:ph-icp}), moves the cloud and repeats until the error stops falling.

\begin{table}[!htbp]\centering\footnotesize\setlength{\tabcolsep}{3pt}\caption{Registration of two views of 400 points rotated by $60^\circ$ about the vertical axis, with Gaussian noise in the points (SVD solution of Eq.~\ref{eq:ph-icp}).}
\begin{tabular}{rrrr}\toprule
\hd{noise $\sigma$ (mm)} & \hd{recovered angle ($^\circ$)} & \hd{error ($^\circ$)} & \hd{residual RMS (mm)}\\\midrule
0.0 & 60.000 & 0.000 & 0.00\\
0.5 & 59.946 & 0.054 & 0.87\\
1.0 & 60.007 & 0.007 & 1.68\\
2.0 & 60.042 & 0.042 & 3.53\\
4.0 & 60.416 & 0.416 & 7.01\\\bottomrule\end{tabular}\end{table}

\exbox{Registration error}{Two views of a plant have 400 matched points spread over about 50 mm in width and 80 mm in height. The scanner noise is $\sigma=1$ mm. Estimate the error in the registration angle and say how many views are needed to close the circle.}\sol \givenline{$\sigma=1$ mm}\ans Six views at $60^\circ$ close the circle ($6\times60^\circ=360^\circ$) \ans With the numbers of the table, at $\sigma=1.0$ mm the angle of one pair is recovered with an error of order $0.01^\circ$ and a residual RMS of about 1.7 mm, which is the noise itself, so the registration adds no bias \ans The error of a chain of five joints adds as the square root of the number of joints if the errors are independent: $\sqrt5\times0.01\approx0.02^\circ$, which at a radius of 25 mm moves a point by $25\times0.02\times\pi/180\approx0.009$ mm, negligible against the 0.6-mm voxel \finans

\section{Measuring volume from the cloud}
The joined cloud is voxelised with a cube of side $v$ (default 0.6 mm) and the volume is the number of occupied voxels times $v^3$; the crown width and the height are the extents of the cloud; the leaf area is the sum of the areas of the triangles of the surface of each leaf, after the points of a leaf have been clustered by their angle around the stem (Zhu et al.\ 2018). The voxel size is a compromise between detail and noise. For a solid body the occupied volume converges to the true volume as $v\to0$; for a surface cloud, the voxels along the surface are the only ones filled, so the estimate is biased unless the interior is filled.

\begin{equation}\boxed{V_{vox}=N_{occ}\,v^{3},\qquad \varepsilon=\frac{V_{vox}-V_{true}}{V_{true}}\approx\frac{3}{2}\frac{v}{r}\ \ (\text{ball of radius }r\text{, surface voxels only}),}\label{eq:ph-vox}\end{equation}
\par\noindent where $V_{vox}$ is the volume estimated (mm$^3$), $N_{occ}$ the number of occupied voxels, $v$ the voxel side (mm), $V_{true}$ the true volume, $\varepsilon$ the relative error and $r$ the radius of the body (mm); for a cloud that covers only the surface the occupied shell has a thickness of about one voxel.



\begin{table}[!htbp]\centering\footnotesize\setlength{\tabcolsep}{3pt}\caption{Voxel volume of a ball of radius 30 mm ($V_{true}=113097$ mm$^3$) with the interior filled, and with the surface only.}
\begin{tabular}{rrrrr}\toprule
\hd{voxel $v$ (mm)} & \hd{solid fill $V$ (mm$^3$)} & \hd{error (\%)} & \hd{surface shell $V$ (mm$^3$)} & \hd{error (\%)}\\\midrule
0.2 & 113100 & 0.00 & 2254 & -98.0\\
0.4 & 113082 & -0.01 & 4449 & -96.1\\
0.6 & 113007 & -0.08 & 6677 & -94.1\\
1.0 & 113081 & -0.01 & 11138 & -90.2\\
1.5 & 112728 & -0.33 & 15964 & -85.9\\
2.5 & 111766 & -1.18 & 24656 & -78.2\\\bottomrule\end{tabular}\end{table}

\exbox{}{A canopy of a seedling is the ellipsoid with semi-axes 40, 40 and 60 mm. Compute its true volume, and the volume found by voxels of 0.6 mm and of 2.5 mm if the interior is filled.}\sol \finans

\exbox{Volume of a canopy by voxels}{A canopy of a seedling is the ellipsoid with semi-axes 40, 40 and 60 mm. Compute its true volume and the volume found by voxels of 0.6 mm and of 2.5 mm if the interior is filled.}\sol \par\indent We know that
\[V=\frac43\pi abc\ti{i}\]
Substituting the values in equation \rfi{i},
\[V=\frac43\pi(40)(40)(60)=402124\ \text{mm$^3$}\] \ans Counting the voxels inside the ellipsoid, $v=0.6$ mm gives $402092$ mm$^3$ (-0.01 per cent error) \ans $v=2.5$ mm gives $402656$ mm$^3$ (0.13 per cent error) \ans A coarse voxel is nearly as good for a solid body; the fine voxel matters for thin leaves, whose thickness is smaller than the voxel \finans

\section{Water productivity and stress traits from the scan}
The scan gives a biomass proxy, $B=\rho_vV$, where $\rho_v$ is a volumetric dry density of the tissue; combined with the water used in the water balance, it gives the transpiration efficiency and the water productivity of each genotype or treatment. The comparison of the well-watered and drought treatments gives the \emph{drought index}

\begin{equation}\boxed{DI=1-\frac{V_{dr}}{V_{ww}},\qquad WP_{B}=\frac{B}{ET_a},\qquad \Delta WP=\frac{WP_{dr}-WP_{ww}}{WP_{ww}},}\label{eq:ph-di}\end{equation}
\par\noindent where $DI$ is the drought index (0 for no effect of drought), $V_{dr}$ and $V_{ww}$ the canopy volumes in the drought and well-watered treatments (cm$^3$), $WP_B$ the biomass water productivity (g\,mm$^{-1}$ per plant, or kg\,m$^{-3}$ for a stand), $B$ the biomass (g), $ET_a$ the actual evapotranspiration (mm) and $\Delta WP$ the relative gain of productivity under drought.



\exbox{Drought index and water productivity from a scan}{A scan gives the canopy volume of 520 cm$^3$ per plant in the well-watered treatment ($ET_a=410$ mm) and 330 cm$^3$ in the drought treatment ($ET_a=270$ mm); the dry density of the tissue is $\rho_v=0.045$ g\,cm$^{-3}$. Find the drought index, the biomass and the water productivities.}\sol \givenline{$ET_a=410$ mm ; $ET_a=270$ mm ; $\rho_v=0.045$ g\,cm$^{-3}$}\ans $DI=1-330/520=0.365$ \ans \[B_{ww}=0.045\times520=23.4\ \text{g}\]
\[B_{dr}=0.045\times330=14.8\ \text{g}\] \ans $WP_{ww}=23.4/410=0.0571$ g\,mm$^{-1}$; $WP_{dr}=14.8/270=0.0550$ g\,mm$^{-1}$ \ans $\Delta WP=-3.6$ per cent: the drought plant used 34 per cent less water and lost 37 per cent of its volume, so it used the water more efficiently \finans

}{}
\banner{Phenotyping layers: the deterministic plant and the volumetric scan}\label{ch:pheno}
Layer P1 grows a plant for $N=140$ days on a soil of eight layers of 15 cm, with no random numbers. Its equations are numbered P.1 to P.28 in the order the model uses them each day. Where the application gives the form but not a species constant, the constant below is an \emph{assumed} typical value, and it is labelled.
\section{Development, light and biomass (P.1 to P.5)}

\exbox{Thermal time (P.1)}{Maize ($T_{\rm base}=8$, $T_{\rm opt}=30$ \textdegree C) grows under the (assumed) daily mean temperature $T(d)=18+9\sin(2\pi(d-60)/200)$. Find the degree days $d_g=\max(0,\min(T,T_{\rm opt}+6)-T_{\rm base})$ on days 0, 30, 70, 100 and the accumulated GDD at day 70 and day 139; if flowering needs $\mathrm{GDD}_{\rm flowering}=800$, on which day is it reached?}\sol  Day 0: $T=18+9\sin(-0.6\pi)=9.44$, $d_g=1.44$. Day 30: $T=10.72$, $d_g=2.72$. Day 70: $T=20.78$, $d_g=12.78$. Day 100: $T=26.56$, $d_g=18.56$. Summing: $\mathrm{GDD}(70)=346$ and $\mathrm{GDD}(139)=1536$ degree days. Flowering (800) is reached on day 99; the stage before flowering is $\mathrm{GDD}/800$, which is 0.43 on day 70.


\exbox{Reference evapotranspiration of the day (P.2)}{The reference evapotranspiration of the day (FAO-56, $G=0$) is $ET_0=\dfrac{0.408\Delta R_n+\gamma\frac{900}{T+273}u_2(e_s-e_a)}{\Delta+\gamma(1+0.34u_2)}$. For $T=28$ \textdegree C, $R_n=18$ MJ m$^{-2}$ d$^{-1}$, $u_2=2$ m s$^{-1}$, $\mathrm{RH}=45$ per cent and $\gamma=0.0674$ kPa \textdegree C$^{-1}$, compute it.}\sol  $e_s=0.6108\exp(17.27\times28/265.3)=3.780$ kPa; $e_a=0.45\,e_s=1.701$ kPa; $\Delta=4098\,e_s/(T+237.3)^2=0.2201$ kPa \textdegree C$^{-1}$. Radiation term: $0.408\times0.2201\times18=1.616$. Aerodynamic term: $0.0674\times900/301\times2\times(2.079)=0.838$. Denominator: $0.2201+0.0674(1+0.68)=0.3333$. So $ET_0=7.36$ mm d$^{-1}$ (about 66 per cent from radiation).


\exbox{Growth from light and from water (P.3 to P.5)}{Take $k=0.55$, RUE $=3.0$ g MJ$^{-1}$ PAR (assumed), $f_T=0.90$, $f_N=0.95$, $f_{\rm sal}=1$, PAR $=10$ MJ m$^{-2}$. Compute the light interception $f_i=1-e^{-k\,\mathrm{LAI}}$ and the radiation-limited growth $dW_{\rm rad}=\mathrm{RUE}f_Tf_Nf_{\rm sal}\mathrm{PAR}\,f_i$ for LAI $=0.5,\dots,5$. Then the water-limited growth $dW_{\rm wat}=(k_{te}/\mathrm{VPD})T_{\rm act}$ for $k_{te}=8$ Pa, VPD $=1.8$ kPa, $T_{\rm act}=3.5$ mm.}\sol  For LAI $=3$: $f_i=1-e^{-1.65}=0.808$, so $dW_{\rm rad}=3.0\times0.9\times0.95\times10\times0.808=20.7$ g m$^{-2}$ d$^{-1}$. Table~\ref{tab:rad} gives the others. For the water limit, $k_{te}/\mathrm{VPD}=8\ \mathrm{Pa}/1800\ \mathrm{Pa}=0.00444$ kg dry matter per kg water, and 3.5 mm of transpiration is 3.5 kg m$^{-2}$, so $dW_{\rm wat}=0.00444\times3.5=0.01556$ kg m$^{-2}$ $=15.6$ g m$^{-2}$. Growth is the smaller of the two, $dW=\min(dW_{\rm rad},dW_{\rm wat})$: for LAI $\le2$ light is limiting ($\le17.1$ g m$^{-2}$ d$^{-1}$ at LAI 2), for LAI $\ge3$ the 15.6 g of water limits it.\begin{table}[!htbp]\centering\small\caption{Light interception and radiation-limited growth.}\label{tab:rad}
\begin{tabular}{lrr}\toprule
\hd{LAI} & \hd{$f_i$} & \hd{$dW_{\rm rad}$ (g m$^{-2}$ d$^{-1}$)}\\\midrule
0.5 & 0.240 & 6.17\\
1 & 0.423 & 10.85\\
2 & 0.667 & 17.11\\
3 & 0.808 & 20.72\\
4 & 0.889 & 22.81\\
5 & 0.936 & 24.01\\\bottomrule\end{tabular}\end{table}


\section{Soil water, leaf water and stomata (P.6 to P.10)}
\exbox{Leaf water potential and relative water content (P.8)}{$\psi=\pi_0/\mathrm{RWC}+P(\mathrm{RWC})$ with turgor $P(\mathrm{RWC})=-\pi_0(\mathrm{RWC}-\mathrm{RWC}_{\rm tlp})/(1-\mathrm{RWC}_{\rm tlp})$ above the turgor-loss point and $0$ below it. With $\pi_0=-1.2$ MPa and $\mathrm{RWC}_{\rm tlp}=0.80$ (assumed), tabulate $\psi$ against RWC and solve for the RWC at $\psi_{\rm leaf}=-1.5$ MPa. Also find $\psi_{\rm leaf}=\psi_{\rm root}-rT_{\rm act}/\mathrm{LAI}$ for $\psi_{\rm root}=-0.4$ MPa, $r=0.25$, $T_{\rm act}=4$ mm, LAI $=2$.}\sol  At RWC $=0.9$: $\pi=-1.2/0.9=-1.333$; $P=1.2(0.9-0.8)/0.2=0.6$; $\psi=-0.733$ MPa. Table~\ref{tab:rwc} gives the others. At $\psi=-1.5$ the bisection returns \textbf{RWC} $=0.800$ (exactly the turgor-loss point: there $P=0$ and $\psi=\pi_0/\mathrm{RWC}=-1.2/0.8=-1.5$ MPa). Flow through the plant: $\psi_{\rm leaf}=-0.4-0.25\times4/2=-0.9$ MPa, which corresponds to RWC $=0.878$.\begin{table}[!htbp]\centering\small\caption{Pressure--volume curve of the leaf.}\label{tab:rwc}
\begin{tabular}{lrrr}\toprule
\hd{RWC} & \hd{$\pi_0/\mathrm{RWC}$} & \hd{$P$} & \hd{$\psi$ (MPa)}\\\midrule
1.00 & -1.200 & 1.200 & 0.000\\
0.95 & -1.263 & 0.900 & -0.363\\
0.90 & -1.333 & 0.600 & -0.733\\
0.85 & -1.412 & 0.300 & -1.112\\
0.80 & -1.500 & 0.000 & -1.500\\
0.75 & -1.600 & 0.000 & -1.600\\
0.70 & -1.714 & 0.000 & -1.714\\\bottomrule\end{tabular}\end{table}


\exbox{Stomata (P.9)}{The stomatal opening is $g_s=1/[1+e^{-b(\psi_{\rm leaf}-\psi_{50})}]$. With steepness $b=3$ MPa$^{-1}$ (assumed), compute $g_s$ for $\psi_{\rm leaf}=-0.5,\dots,-2.5$ and three genotypes: $\psi_{50}=-1.5$ (standard), $-1.0$ (a water saver that closes early) and $-2.0$ (a spender).}\sol  For the standard genotype at $-1.0$ MPa: $g_s=1/(1+e^{-3(-1.0+1.5)})=1/(1+e^{-1.5})=0.818$. At $\psi_{50}$ the opening is exactly 0.5. Table~\ref{tab:gs} lists all; at $-1.5$ MPa the water saver has $g_s=1/(1+e^{1.5})=0.182$ against $0.5$ for the standard genotype, so its transpiration demand is cut to $36$ per cent of the standard one at that potential, which saves water but also growth.\begin{table}[!htbp]\centering\small\caption{Stomatal opening of three genotypes.}\label{tab:gs}
\begin{tabular}{lrrr}\toprule
\hd{$\psi_{\rm leaf}$} & \hd{$\psi_{50}=-1.5$} & \hd{$\psi_{50}=-1.0$} & \hd{$\psi_{50}=-2.0$}\\\midrule
-0.5 & 0.953 & 0.818 & 0.989\\
-1.0 & 0.818 & 0.500 & 0.953\\
-1.5 & 0.500 & 0.182 & 0.818\\
-2.0 & 0.182 & 0.047 & 0.500\\
-2.5 & 0.047 & 0.011 & 0.182\\\bottomrule\end{tabular}\end{table}


\exbox{Expansion growth (P.10)}{Expansion is $\frac1V\frac{dV}{dt}=\varphi\max(0,P-Y)$. With $\varphi=0.4$ d$^{-1}$ MPa$^{-1}$, yield threshold $Y=0.3$ MPa (assumed) and the turgor of the pressure--volume curve above (), find the relative growth rate and the leaf-area doubling time.}\sol  RWC $0.95$: $P=1.2(0.15)/0.2=0.9$, rate $=0.4(0.9-0.3)=0.24$ d$^{-1}$, doubling time $\ln2/0.24=2.89$ d. RWC $0.90$: $0.4(0.6-0.3)=0.12$ d$^{-1}$, $5.78$ d. RWC $0.85$: $P=0.3=Y$, rate $0$: growth stops, although the leaf still holds $85$ per cent of its water and the stomata are partly open. Because turgor falls before photosynthesis does, area and height slow first.


\exbox{Canopy temperature and CWSI (P.11)}{The canopy--air temperature difference is $T_c-T_a=\dfrac{r_aR_n}{\rho c_p}\dfrac{\gamma^*}{\Delta+\gamma^*}-\dfrac{\mathrm{VPD}}{\Delta+\gamma^*}$ with $\gamma^*=\gamma(1+r_c/r_a)$, and CWSI $=(dT-dT_{\rm low})/(dT_{\rm up}-dT_{\rm low})$. Take $T_a=30$ \textdegree C ($\Delta=0.2434$ kPa K$^{-1}$), $R_n=400$ W m$^{-2}$, $r_a=40$ s m$^{-1}$, $\rho c_p=1200$ J m$^{-3}$ K$^{-1}$, $\mathrm{VPD}=2.0$ kPa, $\gamma=0.0674$. Find $T_c-T_a$ for canopy resistances $r_c=0,50,100,200,500$ s m$^{-1}$ and the CWSI.}\sol  Limits: a transpiring canopy ($r_c=0$, $\gamma^*=\gamma$): $dT_{\rm low}=-3.54$ K; a non-transpiring one ($r_c\to\infty$, $\gamma^*\to\infty$): $dT_{\rm up}=r_aR_n/(\rho c_p)=13.33$ K. For $r_c=100$: $\gamma^*=0.0674\times3.5=0.2359$, so $dT=13.33\times0.2359/(0.2434+0.2359)-2.0/(0.2434+0.2359)=2.39$ K. Table~\ref{tab:cwsi} lists the others. The canopy temperature depression is $T_a-T_c=-dT$, positive for a canopy cooler than the air.\begin{table}[!htbp]\centering\small\caption{Canopy temperature difference and stress index.}\label{tab:cwsi}
\begin{tabular}{lrr}\toprule
\hd{$r_c$ (s m$^{-1}$)} & \hd{$T_c-T_a$ (K)} & \hd{CWSI}\\\midrule
0 & -3.54 & 0.000\\
50 & 0.06 & 0.213\\
100 & 2.39 & 0.352\\
200 & 5.24 & 0.520\\
500 & 8.79 & 0.731\\\bottomrule\end{tabular}\end{table}


\section{Carbon, roots and soil strength (P.12 to P.14)}
\exbox{Carbon to the roots (P.12)}{The share of new carbon that goes to the roots is $f_r=\mathrm{clip}\bigl[(f_{r0}+(f_{rA}-f_{r0})\,\mathrm{stage})\,c_{\rm root}(1+0.9w_s),\,0.03,\,0.7\bigr]$ with $w_s=1-T_{\rm act}/T_p$. Wheat: $f_{r0}=0.42$, $f_{rA}=0.07$. Take $c_{\rm root}=1$ and find $f_r$ at stage $0,0.5,1$ for $T_{\rm act}/T_p=1,0.75,0.5$.}\sol  Stage $0.5$ and $T_{\rm act}/T_p=0.75$: base $0.42+(0.07-0.42)(0.5)=0.245$; stress $w_s=0.25$, factor $1+0.225=1.225$; so $f_r=0.245\times1.225=0.300$. The results are in Table~\ref{tab:fr}. Water shortage raises the root share (up to 45 per cent more at $T_{\rm act}/T_p=0.5$) and the stage lowers it: functional equilibrium.\begin{table}[!htbp]\centering\small\caption{Carbon share to roots.}\label{tab:fr}
\begin{tabular}{lrrr}\toprule
\hd{Stage} & \hd{$T_{\rm act}/T_p$} & \hd{$w_s$} & \hd{$f_r$}\\\midrule
0 & 1.00 & 0.00 & 0.420\\
0 & 0.75 & 0.25 & 0.515\\
0 & 0.50 & 0.50 & 0.609\\
0.5 & 1.00 & 0.00 & 0.245\\
0.5 & 0.75 & 0.25 & 0.300\\
0.5 & 0.50 & 0.50 & 0.355\\
1.0 & 1.00 & 0.00 & 0.070\\
1.0 & 0.75 & 0.25 & 0.086\\
1.0 & 0.50 & 0.50 & 0.102\\\bottomrule\end{tabular}\end{table}


\exbox{Soil strength and root elongation (P.14)}{Soil strength is $Q=0.5\,e^{4.5(\rho_b-1.3)}[1+(1-S)]$ (MPa) and the relative elongation is $e/e_0=1/[1+(Q/2.5)^3]$. Compute both for bulk densities $1.2$ to $1.7$ g cm$^{-3}$ at relative water content $S=0.9$ and $0.6$, and find the bulk density that halves the elongation at $S=0.6$.}\sol  For $\rho_b=1.5$, $S=0.6$: $Q=0.5e^{0.9}(1.4)=0.5(2.4596)(1.4)=1.722$ MPa; $e/e_0=1/(1+(0.6887)^3)=1/(1.3267)=0.754$. Table~\ref{tab:Q} lists the rest. Halving requires $Q=2.5$: $0.5e^{4.5(\rho_b-1.3)}(1.4)=2.5$, $e^{4.5(\rho_b-1.3)}=3.571$, $\rho_b-1.3=\ln3.571/4.5=0.2831$, so $\rho_b=1.583$ g cm$^{-3}$. Drying the soil to $S=0.6$ from $0.9$ raises $Q$ by 27 per cent (factor $1.4/1.1$), as strength rises when the soil dries.\begin{table}[!htbp]\centering\small\caption{Soil strength and root elongation.}\label{tab:Q}
\begin{tabular}{lrrr}\toprule
\hd{$\rho_b$} & \hd{$S$} & \hd{$Q$ (MPa)} & \hd{$e/e_0$}\\\midrule
1.2 & 0.9 & 0.351 & 0.997\\
1.2 & 0.6 & 0.446 & 0.994\\
1.3 & 0.9 & 0.550 & 0.989\\
1.3 & 0.6 & 0.700 & 0.979\\
1.4 & 0.9 & 0.863 & 0.961\\
1.4 & 0.6 & 1.098 & 0.922\\
1.5 & 0.9 & 1.353 & 0.863\\
1.5 & 0.6 & 1.722 & 0.754\\
1.6 & 0.9 & 2.122 & 0.621\\
1.6 & 0.6 & 2.700 & 0.442\\
1.7 & 0.9 & 3.327 & 0.298\\
1.7 & 0.6 & 4.235 & 0.171\\\bottomrule\end{tabular}\end{table}


\exbox{Root elongation (P.13)}{Root tip elongation is $\dfrac{dl}{dt}=e_0\,2^{(T_s-20)/10}\Bigl(\dfrac e{e_0}\Bigr)f_wf_{\rm age}$ with $f_w=\mathrm{clip}\bigl((\theta_{\rm rel}-0.06)/0.5;0.05,1\bigr)$. For an axis with $e_0=2.5$ cm d$^{-1}$ at $T_s=16$ \textdegree C, $\theta_{\rm rel}=0.5$, $e/e_0=0.754$ and $f_{\rm age}=0.8$, find the elongation and the time to reach 60 cm.}\sol  Temperature factor $2^{-0.4}=0.758$; $f_w=(0.5-0.06)/0.5=0.88$. Elongation: $2.5\times0.758\times0.754\times0.88\times0.8=1.006$ cm d$^{-1}$, so a tip needs about 60 days to reach 60 cm at these conditions; in the wetter, looser soil ($e/e_0=1$, $f_w=1$, $T_s=20$) the same axis would need $60/(2.5\times0.8)=30$ days.


\section{Irrigation, mulch, nitrogen and salt in the plant (P.15 to P.19)}
\exbox{Project irrigation event (P.16)}{Irrigation (P.16): if the sensor suction $h_s\ge h_{\rm trig}$ and the minimum interval has passed, the root zone is refilled to field capacity, $d_{\rm net}=\min(\sum(\theta_{fc}-\theta_z)\Delta z,D_{\max})$, and the gross depth is $I=d_{\rm net}/E_a$. The sprinkler has minimum interval 2 d, $D_{\max}=35$ mm, $E_a=0.75$; surface 3 d, $60$ mm, $0.60$; drip 1 d, 10 mm, $0.90$. Eight layers of $\Delta z=15$ cm have deficits $\theta_{fc}-\theta_z=0.02,0.03,0.04,0.05,0.05,0.04,0.03,0.02$. Find the net and gross depth for each method.}\sol  The total deficit is $\sum(\theta_{fc}-\theta_z)\Delta z=15(0.02+0.03+0.04+0.05+0.05+0.04+0.03+0.02)=15(0.28)=4.2$ cm $=42$ mm. Sprinkler: $d_{\rm net}=\min(42,35)=35$ mm, $I=35/0.75=46.7$ mm. Surface: $\min(42,60)=42$ mm, $I=42/0.60=70$ mm. Drip: $\min(42,10)=10$ mm, $I=10/0.90=11.1$ mm per event (and the drip line waters only when the layer at the emitter depth is at least 12 per cent depleted). The surface method needs 70 mm per event against 46.7 for the sprinkler for the same deficit; a drip system needs about four small events to replace the same deficit.


\exbox{Mulch (P.17)}{Mulch (P.17) gives $E_p'=f_mE_p$ and $P'=\max(0,P-C_m)$. For $f_m=0.4$, $C_m=2$ mm and a day with $E_p=4.5$ mm and a rain event of $P=1.5$, $4$ and $12$ mm, find the soil evaporation and the effective rain.}\sol  $E_p'=0.4\times4.5=1.8$ mm. $P'=\max(0,1.5-2)=0$ for the 1.5 mm event (the mulch keeps all of it), $2.0$ mm for the 4 mm event, $10.0$ mm for the 12 mm. Over a season with many small events the interception can be substantial: if the rain of 90 events of mean depth 1.5 mm (135 mm) never exceeds $C_m$, none of it reaches the soil.


\exbox{Salt in the root zone (P.19)}{The root-zone solution of volume $V$ (twice the plant-available water of the column) follows $V\,dC/dt=I\,EC_w+0.05P-DC$, $EC_e=C/2$, $f_{\rm sal}=1-\frac s{100}\max(0,EC_e-EC_{\rm th})$, $\psi_{\rm soil}\leftarrow\psi_{\rm soil}-0.036EC_e$. With $V=200$ mm (assumed), a daily irrigation of $I=4$ mm of $EC_w=2$ dS m$^{-1}$, no rain, a drainage of $D=1.5$ mm d$^{-1}$ and an initial $C=1$, integrate by daily steps and find the steady state; for wheat ($EC_{\rm th}=6$, $s=7.1$) compute $f_{\rm sal}$.}\sol  One step: $C_1=1+(4\times2-1.5\times1)/200=1.0325$. The steady state sets the right-hand side to zero: $C^*=I\,EC_w/D=8/1.5=5.33$ dS m$^{-1}$, so $EC_e=2.67$. The relaxation time is $V/D=200/1.5=133$ days, so after 60 days $C=2.57$ (the analytic $C(t)=C^*-(C^*-C_0)e^{-Dt/V}=1.31$ agrees). At $C^*$, $EC_e=2.67$ is below the threshold 6, so $f_{\rm sal}=1$ and the osmotic term is $-0.036\times2.67=-0.096$ MPa. If the drainage is cut to $0.3$ mm d$^{-1}$: $C^*=26.7$, $EC_e=13.3$, $f_{\rm sal}=1-0.071(13.3-6)=0.48$, a loss of 52 per cent.


\section{Fit to the Daisy data and genotype call (P.20, P.21)}
\exbox{Objective of the fit (P.20)}{The plant is compared with the Daisy runs by the objective $J(\theta)=\sum_k\bigl((m_k(\theta)-r_k)/\tau_k\bigr)^2$ over the relative yield, the transpiration share $\tau=T/(T+E_s)$, the mulch effect and five points of the E2 curve. The multipliers $\theta=(f_{\rm RUE},f_{kte},f_{qG})$ are searched one at a time with a step that is halved when no move helps. Using the assumed model values $m_k$, Daisy values $r_k$ and tolerances $\tau_k$ of Table~\ref{tab:J}, compute $J$.}\sol  Each term is $((m-r)/\tau)^2$. For the relative yield: $(0.86-0.90)/0.05=-0.8$, squared $0.64$. Summing the eight terms gives $J=4.447$. A one-parameter search tries $f_{\rm RUE}\to1.1f_{\rm RUE}$ and keeps the change if $J$ falls; if neither $\pm$ step helps the step is halved, within the bounds 0.6 to 1.6 (RUE), 0.5 to 2.2 ($k_{te}$) and 0.25 to 4 ($q_G$), in at most 45 model runs.\begin{table}[!htbp]\centering\small\caption{Objective of the fit to the Daisy data.}\label{tab:J}
\begin{tabular}{lrrrr}\toprule
\hd{Target} & \hd{Model $m$} & \hd{Daisy $r$} & \hd{$\tau$} & \hd{$((m-r)/\tau)^2$}\\\midrule
relative yield & 0.86 & 0.90 & 0.05 & 0.640\\
transpiration share & 0.52 & 0.55 & 0.05 & 0.360\\
mulch effect on $E$ & 0.31 & 0.28 & 0.06 & 0.250\\
E2 point 1 & 0.97 & 1.00 & 0.04 & 0.563\\
E2 point 2 & 0.93 & 0.95 & 0.04 & 0.250\\
E2 point 3 & 0.88 & 0.90 & 0.04 & 0.250\\
E2 point 4 & 0.74 & 0.80 & 0.05 & 1.440\\
E2 point 5 & 0.55 & 0.50 & 0.06 & 0.694\\\bottomrule\end{tabular}\end{table}
\textbf{Reading.} The largest contributions are the points where the model departs by more than its tolerance; the sum is the quantity that the search minimises. A value of $J$ near the number of targets (8) means the model is off by about one tolerance on average.


\exbox{Genotype probabilities (P.21)}{The genotype call weights each genotype by $p_g\propto\exp\bigl(-(d_g-d_{\min})/(0.03+0.25d_{\min})\bigr)$, with $d_g$ the mean squared difference between its traits and the observed traits in units of the trait spread. Five genotypes have distances $0.40,0.55,0.20,0.90,1.60$ (assumed). Find the probabilities.}\sol  $d_{\min}=0.20$, so the temperature is $0.03+0.25(0.20)=0.08$. The weights are $e^{-(d-0.20)/0.08}$: for $d=0.40$, $e^{-2.5}=0.0821$; for $d=0.55$, $e^{-4.375}=0.0126$; $d=0.20$ gives 1; $d=0.90$ gives $e^{-8.75}=1.6\times10^{-4}$; $d=1.60$ gives $2.8\times10^{-7}$. The sum is 1.0948, so the probabilities are in Table~\ref{tab:geno}.\begin{table}[!htbp]\centering\small\caption{Genotype call.}\label{tab:geno}
\begin{tabular}{lrrrr}\toprule
\hd{Genotype} & \hd{$d_g$} & \hd{$d_g-d_{\min}$} & \hd{Weight} & \hd{$p_g$}\\\midrule
G1 & 0.40 & 0.20 & 0.0821 & 0.075\\
G2 & 0.55 & 0.35 & 0.0126 & 0.011\\
G3 & 0.20 & 0.00 & 1.0000 & 0.913\\
G4 & 0.90 & 0.70 & 0.0002 & 0.000\\
G5 & 1.60 & 1.40 & 0.0000 & 0.000\\\bottomrule\end{tabular}\end{table}
The call is confident (about 91 per cent for G3) because the best distance is much smaller than the next; had the distances been $0.20$ and $0.22$ the weights would be $1$ and $0.78$ and the call would be uncertain. The temperature grows with $d_{\min}$, so a poor best fit never gives a sharp call.


\exbox{Solidity, DSI and STI (P.22)}{The projected shoot area of a plant is $A_{\rm proj}=420$ cm$^2$ and the convex hull around the leaves is $A_{\rm hull}=560$ cm$^2$. Four genotypes have grain under stress $Y_s$ and without stress $Y_p$ (Mg ha$^{-1}$, assumed) as in Table~\ref{tab:dsi}. Compute the solidity, the drought susceptibility index DSI $=(1-Y_s/Y_p)/(1-\bar Y_s/\bar Y_p)$ and the stress tolerance index STI $=Y_sY_p/\bar Y_p^2$.}\sol  Solidity $=A_{\rm proj}/A_{\rm hull}=0.750$. The means are $\bar Y_s=4.500$, $\bar Y_p=6.500$, so $1-\bar Y_s/\bar Y_p=0.3077$. For G2: $1-4.2/7.1=0.4085$, DSI $=0.4085/0.3077=1.327$; STI $=4.2\times7.1/6.500^2=0.706$.\begin{table}[!htbp]\centering\small\caption{Drought susceptibility and stress tolerance indices.}\label{tab:dsi}
\begin{tabular}{lrrrr}\toprule
\hd{Genotype} & \hd{$Y_s$} & \hd{$Y_p$} & \hd{DSI} & \hd{STI}\\\midrule
G1 & 5.1 & 6.8 & 0.812 & 0.821\\
G2 & 4.2 & 7.1 & 1.327 & 0.706\\
G3 & 5.6 & 6.2 & 0.315 & 0.822\\
G4 & 3.1 & 5.9 & 1.542 & 0.433\\\bottomrule\end{tabular}\end{table}
\textbf{Reading.} DSI above 1 means the genotype loses more than the mean under stress (G2, G4); below 1 it is more tolerant (G3). STI rewards genotypes that are high in both environments: here G1 and G3 have the highest values.


\exbox{Depth of 95 per cent of the roots (P.23)}{Root length (cm) in the eight layers of 15 cm is $420,300,210,140,90,60,35,15$ (assumed). Find the depth $D_{95}$ at which the cumulative length reaches 95 per cent of the total, the root length density RLD $=L/(A\Delta z)$ of the top layer for a ground area $A=900$ cm$^2$, and the root-to-shoot ratio for $W_{\rm root}=18$ g and $W_{\rm shoot}=62$ g.}\sol  Total $=1270$ cm. Cumulative shares: 0.331, 0.567, 0.732, 0.843, 0.913, 0.961, 0.988, 1.000. The 95 per cent mark lies in layer 6 (cumulative 0.913 before it and 0.961 after); by linear interpolation inside the layer, $D_{95}=86.6$ cm. RLD of the top layer $=420/(900\times15)=0.0311$ cm cm$^{-3}$. Root-to-shoot ratio $=18/62=0.290$.


\section{The drought outlook and the stress calendar (P.24 to P.28)}
\exbox{Analogue years (P.25)}{Each of the 20 years of the weather has the rain of the days that have no forecast, $S_i=\sum_{k\in\rm future}P_i(k)$ (mm), $412,388,455,301,517,366,298,470,405,342,380,436,290,501,352,421,330,399,463,377$ (assumed). Find the dry, typical and wet analogue years $i_{\rm dry}=\arg\min S_i$, $i_{\rm wet}=\arg\max S_i$ and $i_{\rm typ}=$ median.}\sol  The minimum is $290$ (year 13), the maximum 517 (year 5). The sorted values are 290, 298, 301, 330, 342, 352, 366, 377, 380, 388, 399, 405, 412, 421, 436, 455, 463, 470, 501, 517; with 20 values the median is the mean of the 10th and 11th, $(388+399)/2=393.5$ mm, and the typical year is the one nearest to it: year 2 ($S=388$). The three futures are real years of the record, so each has its own dry spells in the order in which they happened; they are not a forecast.


\exbox{Periods of drought stress (P.26, P.27)}{The root-zone water is $a=\sum f_i\theta_i/\sum f_i$, the FAO-56 stress coefficient is $K_s=\min(1,a/(1-p))$ and a day is a stress day if $a<1-p$ ($p=0.55$ for maize and wheat). The root-zone values of 16 days are in Table~\ref{tab:ks}. Find the stress days, join gaps of at most 2 days, and report the periods (kept if $n_j\ge2$) with $a_{\min}$ and mean $K_s$.}\sol  The threshold is $1-p=0.45$. Stress days: those with $a<0.45$. The raw runs are [4,6], [8,10], [15,16]. Gaps of at most 2 days are filled, which merges these into [4,10], [15,16]. Periods kept: days 4 to 10 ($n=7$, $a_{\min}=0.33$, mean $K_s=0.867$); days 15 to 16 ($n=2$, $a_{\min}=0.35$, mean $K_s=0.822$).\begin{table}[!htbp]\centering\small\caption{Root-zone water and the stress coefficient.}\label{tab:ks}
\begin{tabular}{lrrr}\toprule
\hd{Day} & \hd{$a$} & \hd{$K_s$} & \hd{Stress}\\\midrule
1 & 0.62 & 1.000 & -\\
2 & 0.55 & 1.000 & -\\
3 & 0.48 & 1.000 & -\\
4 & 0.40 & 0.889 & stress\\
5 & 0.38 & 0.844 & stress\\
6 & 0.44 & 0.978 & stress\\
7 & 0.52 & 1.000 & -\\
8 & 0.36 & 0.800 & stress\\
9 & 0.33 & 0.733 & stress\\
10 & 0.37 & 0.822 & stress\\
11 & 0.50 & 1.000 & -\\
12 & 0.60 & 1.000 & -\\
13 & 0.58 & 1.000 & -\\
14 & 0.46 & 1.000 & -\\
15 & 0.39 & 0.867 & stress\\
16 & 0.35 & 0.778 & stress\\\bottomrule\end{tabular}\end{table}


\exbox{Risk class and what the roots suffer (P.26, P.28)}{The risk class is set by $R=Y_{\rm rainfed}/Y_{\rm irrigated}$: low if $R\ge0.90$, moderate if $0.75\le R<0.90$, high if $0.50\le R<0.75$, severe below 0.50. For a dry future the rainfed grain is $3.1$ Mg ha$^{-1}$ and the irrigated $6.4$; typical $5.6/6.4$; wet $6.2/6.4$. Classify them, and compute the root measures of one period (P.28): root growth $\Delta L_k=(0.8,0.7,0.5,0.4)$ cm d$^{-1}\times$ plant against the bench plant $(1.0,1.0,1.0,0.9)$, deepest root $D=62$ cm against bench $70$ cm, and root lengths $L_i=(120,90,70,40,20)$ cm in layers with $\theta_i=(0.7,0.55,0.45,0.30,0.25)$.}\sol  $R=3.1/6.4=0.484$: severe; $5.6/6.4=0.875$: moderate; $6.2/6.4=0.969$: low. P.28: $G=100(0.8+0.7+0.5+0.4)/(1.0+1.0+1.0+0.9)=100(2.4)/(3.9)=61.5$ per cent (the roots grew about 62 per cent as much as the watered plant); $\Delta D=70-62=8$ cm shallower; layers with $\theta<0.5$ are layers 3, 4, 5: $F=100(70+40+20)/(120+90+70+40+20)=100(130)/340=38.2$ per cent of the root length lies in layers holding less than half of their plant-available water.


\section{Volumetric phenotyping (P2): rebuilding the scan}
Layer P2 measures the plant of P1 as a three-dimensional scanner would: views from a turning camera, calibration spheres, registration, a voxel cloud and traits.

\exbox{Joining two views by the spheres}{A calibration ring of six spheres of centres at radius $0.35$ m and angles $0,60,\dots,300$\textdegree\ is seen by view 2 turned by $60$\textdegree\ and shifted by $(0.020,-0.010)$ m, with centre-fit noise of $1$ mm. Four spheres (those at $0,60,120,180$\textdegree) are seen by both views. Estimate the rigid transform from view 2 back to view 1 (Procrustes/Kabsch) and its error.}\sol  1. Centroids: of the four observed centres $c_A=(-0.1109,0.0658)$ m, of the reference centres $c_B=(0.0000,0.1516)$ m. 2. Cross-covariance $H=\sum(A_i-c_A)(B_i-c_B)^T$; its singular value decomposition $H=U\Sigma V^T$ gives the rotation $R=V\,\mathrm{diag}(1,\det(VU^T))U^T$, with $\det(VU^T)=+1$ here so that $R$ is a proper rotation. 3. The rotation angle is $-60.07$\textdegree\ (truth: $-60$\textdegree\ to undo the turn), an error of $0.067$\textdegree. 4. Translation $t=c_B-Rc_A=(-0.0017,0.0226)$ m. 5. The residual after alignment has RMS $0.89$ mm, about the centre-fit noise of 1 mm. At the crown radius of $0.4$ m the angle error moves a point by $0.4\times0.00116=0.47$ mm.


\exbox{Accumulated registration error}{The transforms are chained to view 1: $R_0,\ R_0R_1,\ R_0R_1R_2,\dots$, so that the error of each link adds. If every link has a turn error of $\delta=0.3$\textdegree\ (a calibrated rig), compute the accumulated angle and the displacement at a point $0.4$ m from the axis for views $1$ to $6$; compare with the case of independent random errors that add in quadrature.}\sol  If the errors are all of the same sign the angle after $k$ links is $k\delta$ and the displacement is $r\,k\delta$ (in radians). For $k=6$: $1.8$\textdegree\ $=0.0314$ rad, displacement $0.4\times0.0314=12.6$ mm. If the link errors are independent with standard deviation $\delta$, the accumulated standard deviation is $\sqrt k\,\delta$, $0.73$\textdegree, so $5.1$ mm at $k=6$ (Table~\ref{tab:chain}). The last views drift most; the ICP refinement on the joined cloud is there to reduce this.\begin{table}[!htbp]\centering\small\caption{Accumulated registration error at 0.4 m.}\label{tab:chain}
\begin{tabular}{lrrr}\toprule
\hd{View $k$} & \hd{Worst case (deg)} & \hd{Worst case (mm)} & \hd{Random (mm)}\\\midrule
1 & 0.30 & 2.09 & 2.09\\
2 & 0.60 & 4.19 & 2.96\\
3 & 0.90 & 6.28 & 3.63\\
4 & 1.20 & 8.38 & 4.19\\
5 & 1.50 & 10.47 & 4.68\\
6 & 1.80 & 12.57 & 5.13\\\bottomrule\end{tabular}\end{table}


\exbox{Height, crown, volume and leaf area from a voxel cloud}{A plant is modelled as a prolate ellipsoid with semi-axes $a=b=0.15$ m and $c=0.30$ m, standing on the ground (centre at $z=0.30$ m). A scan of 60,000 points is placed in voxels of $v=0.02$ m. The application measures: height = 99.8th percentile of $z$; crown width = longest horizontal span of the convex hull; canopy volume = occupied cubes; leaf area = surface voxels $\times v^2/1.5$. Compute them and compare with the truth.}\sol  Height: the 99.8th percentile of the point heights is 0.599 m (truth $0.600$ m; the percentile discards the top 0.2 per cent of points, so it is slightly below the apex). Crown width: the hull of the points projected on the ground has longest span 0.300 m (truth $2a=0.300$ m). Volume: the true volume is $\tfrac43\pi abc=28.27$ litres (0.02827 m$^3$); counting the cubes of $0.02$ m whose centres lie inside gives 3528 cubes, $V=28.22$ litres, an error of -0.2 per cent. Leaf area: the cloud occupies 1648 surface voxels, so $LA=1648\times0.02^2/1.5=0.439$ m$^2$ against the true surface of the ellipsoid 0.483 m$^2$ (Knud Thomsen approximation), a ratio of 0.91. The factor 1.5 is the application's convention; the voxel count depends on the voxel size (a coarse voxel counts the thickness of the leaf), as the application states. The leaf area index for a ground area of $0.25$ m$^2$ is $LAI=LA/A=1.76$.


\exbox{Leaf area against the voxel size}{Repeat the leaf-area measurement for voxel sizes 0.01, 0.02, 0.04 and 0.08 m with the same cloud.}\sol  The number of surface voxels falls roughly as $v^{-2}$ so that $N v^2$ would be constant for a thin surface; a coarser voxel also counts the thickness, so the estimate departs from the true area. Table~\ref{tab:vox} shows the area and its ratio to the true surface (0.483 m$^2$).\begin{table}[!htbp]\centering\small\caption{Leaf area against voxel size.}\label{tab:vox}
\begin{tabular}{lrrr}\toprule
\hd{Voxel (m)} & \hd{Surface voxels} & \hd{$LA$ (m$^2$)} & \hd{$LA$/true}\\\midrule
0.01 & 6194 & 0.413 & 0.85\\
0.02 & 1648 & 0.439 & 0.91\\
0.04 & 420 & 0.448 & 0.93\\
0.08 & 96 & 0.410 & 0.85\\\bottomrule\end{tabular}\end{table}
\textbf{Reading.} The estimate is not corrected for the voxel size in the application, and the error shown there is honest about it: a scanner's resolution is part of the measurement.


\exbox{Volumetric stress index}{The volumetric stress index is $\mathrm{VSI}=V_{\rm drought}/V_{\rm project}$. For a drought plant with canopy volume $18.4$ litres and a project plant with $26.9$ litres (assumed), and for scans every 10 days with volumes (drought, project) of $(2.1,2.3),(5.8,7.2),(11.9,15.8),(18.4,26.9)$ litres, compute the VSI each time and the transpiration per litre of canopy if the drought plant transpired 9.6 L and the project plant 14.2 L by the last date.}\sol  $\mathrm{VSI}=2.1/2.3=0.913$; $5.8/7.2=0.806$; $11.9/15.8=0.753$; $18.4/26.9=0.684$. The drought plant ends 32 per cent smaller. Transpiration per litre of canopy: drought $9.6/18.4=0.522$ L per litre; project $14.2/26.9=0.528$ L per litre; the two are within 1 per cent, so the drought plant is smaller in proportion to the water it used: it did not lose efficiency, it lost size.
\IfFileExists{extra/app10.tex}{\subsection{The output of the application, panel by panel}
The phenotyping layers have two pictures: the deterministic plant of the model and the volumetric scan of a real plant by a camera that turns round it.

The layer draws the plant of the model for the chosen crop and day, with the height, the leaf area index, the roots and the cover.

\appfig{dph_0}{0.96}{Deterministic phenotyping: the plant of the model through the season.}

After the run the numbers of the plant (height, LAI, biomass) are given for the days of the season.

\appfig{dph_1}{0.96}{Run simulation: the phenotype for the settings of the user.}

The page is a guided tour of the six steps from the rig to the measured volume.

\appfig{vph_0}{0.96}{Volumetric phenotyping: the plant, the camera rig and the six steps of the scan.}

The arm turns, a view arrives at each position and the views are joined into one cloud of points.

\appfig{vph_1}{0.96}{Step 2: the camera turns round the plant; the views arrive and are joined.}

The simulation can be played step by step to see the point cloud grow.

\appfig{vph_2}{0.96}{Step 2 of 6 (the camera turns round the plant), with the moving simulation of the scan.}

The same rig scans the plant on seven dates of the season to follow its growth.

\appfig{vph_3}{0.96}{Step 3: growth from emergence to the end of the season.}

The volume of the point cloud (voxel counting) is compared with the model value, an error check of the method.

\appfig{vph_4}{0.96}{Step 5: volume and form measured on the joined cloud against the model for the same day.}

The measured volumes form the growth curve and the line is the deterministic model.

\appfig{vph_5}{0.96}{Step 7: through the season, a scan every ten days (24 of 24 done); the line is the model.}

The last panel states the inputs of the layer and the numbers that depend on it.

\appfig{vph_6}{0.96}{Step 8: what feeds this layer and what it feeds.}

}{}
\IfFileExists{extra/long10.tex}{\subsection{Complete registration and volume of a scan}
\exbox{Registration by SVD and voxel volume}{A box 2 $\times$ 3 $\times$ 4 cm is scanned in a second pose. Recover the rigid motion by the SVD method, check the residual, and find the volume by voxel counting.}\sol
\ans Eight corner points $P_i$ of a box 2 $\times$ 3 $\times$ 4 cm are scanned in a second pose $Q_i=RP_i+t$ with $R$ a rotation of 25$^\circ$ about $z$ and $t=(1.5,-0.7,0.4)$
\ans Centroids: $\bar P=(1.00, 1.50, 2.00)$, $\bar Q=(1.772, 1.082, 2.400)$
\ans Cross-covariance $H=\sum(P_i-\bar P)(Q_i-\bar Q)^{T}$ has singular values $32.00, 18.00, 8.00$
\ans Rotation from the SVD $H=U\Sigma V^T$: $R=VU^T$ gives an angle $\arctan(R_{21}/R_{11})=25.000^\circ$ (true $25^\circ$) and $\det R=1.000$
\ans Translation $t=\bar Q-R\bar P=(1.500, -0.700, 0.400)$
\ans Residual of the registration $\sqrt{\frac{1}{N}\sum|RP_i+t-Q_i|^2}=3.2e-16$ cm
\ans Voxel volume: with voxels of 0.5 cm, the box contains $4\times6\times8=192$ voxels of 0.125 cm$^3$, giving $192\times0.125=24$ cm$^3$ (true 24 cm$^3$) \finans


}{}
\IfFileExists{extra/proj10.tex}{\banner{Design project 10: ranking genotypes in a managed-stress trial}\label{sec:proj10}
\textbf{Brief.} Five genotypes were grown in a stress trial. Their yield, seasonal evapotranspiration, harvest index and canopy temperature at midday are in the table. The task is to rank them by water productivity and stress index and to recommend the genotype for a water-scarce scheme.

\begin{table}[!htbp]\centering\small\caption{Five genotypes in a managed-stress trial.}
\begin{tabular}{lrrrrr}\toprule
\hd{Genotype} & \hd{Yield (Mg/ha)} & \hd{$ET_a$ (mm)} & \hd{WP (kg/m$^3$)} & \hd{HI} & \hd{Canopy temp. (C)}\\\midrule
G1 & 9.8 & 455 & 2.15 & 0.42 & 31.0\\
G2 & 10.6 & 490 & 2.16 & 0.45 & 29.5\\
G3 & 8.9 & 410 & 2.17 & 0.40 & 33.0\\
G4 & 11.2 & 540 & 2.07 & 0.43 & 28.5\\
G5 & 10.1 & 445 & 2.27 & 0.47 & 30.5\\\bottomrule\end{tabular}\end{table}

\begin{table}[!htbp]\centering\small\caption{Crop water stress index from canopy temperature ($T_a=34$ C, lower limit $T_a-3$, upper limit $T_a+4$).}
\begin{tabular}{lrr}\toprule
\hd{Genotype} & \hd{$T_c$ (C)} & \hd{CWSI}\\\midrule
G1 & 31.0 & 0.00\\
G2 & 29.5 & -0.21\\
G3 & 33.0 & 0.29\\
G4 & 28.5 & -0.36\\
G5 & 30.5 & -0.07\\\bottomrule\end{tabular}\end{table}

\textbf{Reading.} The highest water productivity is G5 at 2.27 kg\,m$^{-3}$. The canopy of G3 is the warmest (CWSI 0.29), a sign of the stomata closing early; G4 yields most but uses most water. For a scheme where water is the limit, rank by WP, then confirm it with the CWSI: the genotype with the best WP and a low stress index is the safest.
\fig{g10_geno}{0.95}{Five genotypes of the stress trial. (a) Yield against seasonal evapotranspiration. (b) Water productivity against the crop water stress index.}
}{}
\IfFileExists{extra/ten10.tex}{\subsection{Worked solutions on the ten soils}
A single genotype is grown on every soil: the soil changes the water it can use.

\exbox{Ten-soil example 10a: Coarse loam (S01)}{Coarse loam: genotype G4 with $Y_m=11.5$ Mg\,ha$^{-1}$, $K_y=1.28$ is grown in the 125-day season on $TAW=90$ mm. Find the relative yield, the yield, the water productivity and the biomass if the harvest index is 0.43.}\sol $ET_a/ET_c=747/747=1.000$; $Y/Y_m=1-1.28(1-1.000)=1.000$ \ans $Y=11.50$ Mg\,ha$^{-1}$ \ans $WP=0.15$ kg\,m$^{-3}$ \ans Biomass $B=Y/HI=11.50/0.43=26.7$ Mg\,ha$^{-1}$ \ans

\exbox{Ten-soil example 10a: Coarse sand (S02)}{Coarse sand: genotype G4 with $Y_m=11.5$ Mg\,ha$^{-1}$, $K_y=1.28$ is grown in the 125-day season on $TAW=30$ mm. Find the relative yield, the yield, the water productivity and the biomass if the harvest index is 0.43.}\sol $ET_a/ET_c=747/747=1.000$; $Y/Y_m=1-1.28(1-1.000)=1.000$ \ans $Y=11.50$ Mg\,ha$^{-1}$ \ans $WP=0.15$ kg\,m$^{-3}$ \ans Biomass $B=Y/HI=11.50/0.43=26.7$ Mg\,ha$^{-1}$ \ans

\exbox{Ten-soil example 10a: Sand (S03)}{Sand: genotype G4 with $Y_m=11.5$ Mg\,ha$^{-1}$, $K_y=1.28$ is grown in the 125-day season on $TAW=36$ mm. Find the relative yield, the yield, the water productivity and the biomass if the harvest index is 0.43.}\sol $ET_a/ET_c=747/747=1.000$; $Y/Y_m=1-1.28(1-1.000)=1.000$ \ans $Y=11.50$ Mg\,ha$^{-1}$ \ans $WP=0.15$ kg\,m$^{-3}$ \ans Biomass $B=Y/HI=11.50/0.43=26.7$ Mg\,ha$^{-1}$ \ans

\exbox{Ten-soil example 10a: Loamy sand (S04)}{Loamy sand: genotype G4 with $Y_m=11.5$ Mg\,ha$^{-1}$, $K_y=1.28$ is grown in the 125-day season on $TAW=54$ mm. Find the relative yield, the yield, the water productivity and the biomass if the harvest index is 0.43.}\sol $ET_a/ET_c=747/747=1.000$; $Y/Y_m=1-1.28(1-1.000)=1.000$ \ans $Y=11.50$ Mg\,ha$^{-1}$ \ans $WP=0.15$ kg\,m$^{-3}$ \ans Biomass $B=Y/HI=11.50/0.43=26.7$ Mg\,ha$^{-1}$ \ans

\exbox{Ten-soil example 10a: Sandy loam (S05)}{Sandy loam: genotype G4 with $Y_m=11.5$ Mg\,ha$^{-1}$, $K_y=1.28$ is grown in the 125-day season on $TAW=72$ mm. Find the relative yield, the yield, the water productivity and the biomass if the harvest index is 0.43.}\sol $ET_a/ET_c=747/747=1.000$; $Y/Y_m=1-1.28(1-1.000)=1.000$ \ans $Y=11.50$ Mg\,ha$^{-1}$ \ans $WP=0.15$ kg\,m$^{-3}$ \ans Biomass $B=Y/HI=11.50/0.43=26.7$ Mg\,ha$^{-1}$ \ans

\exbox{Ten-soil example 10a: Loam (S06)}{Loam: genotype G4 with $Y_m=11.5$ Mg\,ha$^{-1}$, $K_y=1.28$ is grown in the 125-day season on $TAW=102$ mm. Find the relative yield, the yield, the water productivity and the biomass if the harvest index is 0.43.}\sol $ET_a/ET_c=747/747=1.000$; $Y/Y_m=1-1.28(1-1.000)=1.000$ \ans $Y=11.50$ Mg\,ha$^{-1}$ \ans $WP=0.15$ kg\,m$^{-3}$ \ans Biomass $B=Y/HI=11.50/0.43=26.7$ Mg\,ha$^{-1}$ \ans

\exbox{Ten-soil example 10a: Silt loam (S07)}{Silt loam: genotype G4 with $Y_m=11.5$ Mg\,ha$^{-1}$, $K_y=1.28$ is grown in the 125-day season on $TAW=120$ mm. Find the relative yield, the yield, the water productivity and the biomass if the harvest index is 0.43.}\sol $ET_a/ET_c=747/747=1.000$; $Y/Y_m=1-1.28(1-1.000)=1.000$ \ans $Y=11.50$ Mg\,ha$^{-1}$ \ans $WP=0.15$ kg\,m$^{-3}$ \ans Biomass $B=Y/HI=11.50/0.43=26.7$ Mg\,ha$^{-1}$ \ans

\exbox{Ten-soil example 10a: Sandy clay loam (S08)}{Sandy clay loam: genotype G4 with $Y_m=11.5$ Mg\,ha$^{-1}$, $K_y=1.28$ is grown in the 125-day season on $TAW=84$ mm. Find the relative yield, the yield, the water productivity and the biomass if the harvest index is 0.43.}\sol $ET_a/ET_c=747/747=1.000$; $Y/Y_m=1-1.28(1-1.000)=1.000$ \ans $Y=11.50$ Mg\,ha$^{-1}$ \ans $WP=0.15$ kg\,m$^{-3}$ \ans Biomass $B=Y/HI=11.50/0.43=26.7$ Mg\,ha$^{-1}$ \ans

\exbox{Ten-soil example 10a: Clay loam (S09)}{Clay loam: genotype G4 with $Y_m=11.5$ Mg\,ha$^{-1}$, $K_y=1.28$ is grown in the 125-day season on $TAW=96$ mm. Find the relative yield, the yield, the water productivity and the biomass if the harvest index is 0.43.}\sol $ET_a/ET_c=747/747=1.000$; $Y/Y_m=1-1.28(1-1.000)=1.000$ \ans $Y=11.50$ Mg\,ha$^{-1}$ \ans $WP=0.15$ kg\,m$^{-3}$ \ans Biomass $B=Y/HI=11.50/0.43=26.7$ Mg\,ha$^{-1}$ \ans

\exbox{Ten-soil example 10a: Silty clay loam (S10)}{Silty clay loam: genotype G4 with $Y_m=11.5$ Mg\,ha$^{-1}$, $K_y=1.28$ is grown in the 125-day season on $TAW=102$ mm. Find the relative yield, the yield, the water productivity and the biomass if the harvest index is 0.43.}\sol $ET_a/ET_c=747/747=1.000$; $Y/Y_m=1-1.28(1-1.000)=1.000$ \ans $Y=11.50$ Mg\,ha$^{-1}$ \ans $WP=0.15$ kg\,m$^{-3}$ \ans Biomass $B=Y/HI=11.50/0.43=26.7$ Mg\,ha$^{-1}$ \ans

}{}
\IfFileExists{extra/tenb10.tex}{\subsection{Worked solutions on ten crops}
The ten crops of the library differ in their demand, their roots and their sensitivity. The same type of question is answered for each of them at $ET_0=5.5$ mm\,d$^{-1}$; the data are typical values of the crops (Doorenbos--Kassam and FAO 56).

\exbox{Crop example 10a: Maize}{Maize: with $Y_m=11.5$ Mg\,ha$^{-1}$ and a seasonal $ET_c\approx660$ mm find the potential water productivity and the harvest index implied by a biomass of 2.4 times the yield.}\sol \givenline{$Y_m=11.5$ Mg\,ha$^{-1}$}\ans \[WP=Y_m/ET_c=11.5\times1000/(660\times10)=1.74\ \text{kg\,m$^{-3}$}\] \ans $B=2.4\times11.5=27.6$ Mg\,ha$^{-1}$; $HI=Y/B=1/2.4=0.42$ \ans $TE=B/T$ with $T\approx0.8ET_c$: \[27.6\times1000/(0.8\times660\times10)=5.23\ \text{kg\,m$^{-3}$}\quad\text{\finans}\]

\exbox{Crop example 10a: Wheat}{Wheat: with $Y_m=8.0$ Mg\,ha$^{-1}$ and a seasonal $ET_c\approx759$ mm find the potential water productivity and the harvest index implied by a biomass of 2.4 times the yield.}\sol \givenline{$Y_m=8.0$ Mg\,ha$^{-1}$}\ans \[WP=Y_m/ET_c=8.0\times1000/(759\times10)=1.05\ \text{kg\,m$^{-3}$}\] \ans $B=2.4\times8.0=19.2$ Mg\,ha$^{-1}$; $HI=Y/B=1/2.4=0.42$ \ans $TE=B/T$ with $T\approx0.8ET_c$: \[19.2\times1000/(0.8\times759\times10)=3.16\ \text{kg\,m$^{-3}$}\quad\text{\finans}\]

\exbox{Crop example 10a: Cotton}{Cotton: with $Y_m=4.5$ Mg\,ha$^{-1}$ and a seasonal $ET_c\approx911$ mm find the potential water productivity and the harvest index implied by a biomass of 2.4 times the yield.}\sol \givenline{$Y_m=4.5$ Mg\,ha$^{-1}$}\ans \[WP=Y_m/ET_c=4.5\times1000/(911\times10)=0.49\ \text{kg\,m$^{-3}$}\] \ans $B=2.4\times4.5=10.8$ Mg\,ha$^{-1}$; $HI=Y/B=1/2.4=0.42$ \ans $TE=B/T$ with $T\approx0.8ET_c$: \[10.8\times1000/(0.8\times911\times10)=1.48\ \text{kg\,m$^{-3}$}\quad\text{\finans}\]

\exbox{Crop example 10a: Tomato}{Tomato: with $Y_m=75.0$ Mg\,ha$^{-1}$ and a seasonal $ET_c\approx683$ mm find the potential water productivity and the harvest index implied by a biomass of 2.4 times the yield.}\sol \givenline{$Y_m=75.0$ Mg\,ha$^{-1}$}\ans \[WP=Y_m/ET_c=75.0\times1000/(683\times10)=10.98\ \text{kg\,m$^{-3}$}\] \ans $B=2.4\times75.0=180.0$ Mg\,ha$^{-1}$; $HI=Y/B=1/2.4=0.42$ \ans $TE=B/T$ with $T\approx0.8ET_c$: \[180.0\times1000/(0.8\times683\times10)=32.94\ \text{kg\,m$^{-3}$}\quad\text{\finans}\]

\exbox{Crop example 10a: Potato}{Potato: with $Y_m=40.0$ Mg\,ha$^{-1}$ and a seasonal $ET_c\approx658$ mm find the potential water productivity and the harvest index implied by a biomass of 2.4 times the yield.}\sol \givenline{$Y_m=40.0$ Mg\,ha$^{-1}$}\ans \[WP=Y_m/ET_c=40.0\times1000/(658\times10)=6.08\ \text{kg\,m$^{-3}$}\] \ans $B=2.4\times40.0=96.0$ Mg\,ha$^{-1}$; $HI=Y/B=1/2.4=0.42$ \ans $TE=B/T$ with $T\approx0.8ET_c$: \[96.0\times1000/(0.8\times658\times10)=18.24\ \text{kg\,m$^{-3}$}\quad\text{\finans}\]

\exbox{Crop example 10a: Soybean}{Soybean: with $Y_m=3.5$ Mg\,ha$^{-1}$ and a seasonal $ET_c\approx658$ mm find the potential water productivity and the harvest index implied by a biomass of 2.4 times the yield.}\sol \givenline{$Y_m=3.5$ Mg\,ha$^{-1}$}\ans \[WP=Y_m/ET_c=3.5\times1000/(658\times10)=0.53\ \text{kg\,m$^{-3}$}\] \ans $B=2.4\times3.5=8.4$ Mg\,ha$^{-1}$; $HI=Y/B=1/2.4=0.42$ \ans $TE=B/T$ with $T\approx0.8ET_c$: \[8.4\times1000/(0.8\times658\times10)=1.60\ \text{kg\,m$^{-3}$}\quad\text{\finans}\]

\exbox{Crop example 10a: Sugar beet}{Sugar beet: with $Y_m=60.0$ Mg\,ha$^{-1}$ and a seasonal $ET_c\approx950$ mm find the potential water productivity and the harvest index implied by a biomass of 2.4 times the yield.}\sol \givenline{$Y_m=60.0$ Mg\,ha$^{-1}$}\ans \[WP=Y_m/ET_c=60.0\times1000/(950\times10)=6.31\ \text{kg\,m$^{-3}$}\] \ans $B=2.4\times60.0=144.0$ Mg\,ha$^{-1}$; $HI=Y/B=1/2.4=0.42$ \ans $TE=B/T$ with $T\approx0.8ET_c$: \[144.0\times1000/(0.8\times950\times10)=18.94\ \text{kg\,m$^{-3}$}\quad\text{\finans}\]

\exbox{Crop example 10a: Sunflower}{Sunflower: with $Y_m=3.0$ Mg\,ha$^{-1}$ and a seasonal $ET_c\approx605$ mm find the potential water productivity and the harvest index implied by a biomass of 2.4 times the yield.}\sol \givenline{$Y_m=3.0$ Mg\,ha$^{-1}$}\ans \[WP=Y_m/ET_c=3.0\times1000/(605\times10)=0.50\ \text{kg\,m$^{-3}$}\] \ans $B=2.4\times3.0=7.2$ Mg\,ha$^{-1}$; $HI=Y/B=1/2.4=0.42$ \ans $TE=B/T$ with $T\approx0.8ET_c$: \[7.2\times1000/(0.8\times605\times10)=1.49\ \text{kg\,m$^{-3}$}\quad\text{\finans}\]

\exbox{Crop example 10a: Onion}{Onion: with $Y_m=50.0$ Mg\,ha$^{-1}$ and a seasonal $ET_c\approx693$ mm find the potential water productivity and the harvest index implied by a biomass of 2.4 times the yield.}\sol \givenline{$Y_m=50.0$ Mg\,ha$^{-1}$}\ans \[WP=Y_m/ET_c=50.0\times1000/(693\times10)=7.22\ \text{kg\,m$^{-3}$}\] \ans $B=2.4\times50.0=120.0$ Mg\,ha$^{-1}$; $HI=Y/B=1/2.4=0.42$ \ans $TE=B/T$ with $T\approx0.8ET_c$: \[120.0\times1000/(0.8\times693\times10)=21.65\ \text{kg\,m$^{-3}$}\quad\text{\finans}\]

\exbox{Crop example 10a: Alfalfa}{Alfalfa: with $Y_m=15.0$ Mg\,ha$^{-1}$ and a seasonal $ET_c\approx506$ mm find the potential water productivity and the harvest index implied by a biomass of 2.4 times the yield.}\sol \givenline{$Y_m=15.0$ Mg\,ha$^{-1}$}\ans \[WP=Y_m/ET_c=15.0\times1000/(506\times10)=2.96\ \text{kg\,m$^{-3}$}\] \ans $B=2.4\times15.0=36.0$ Mg\,ha$^{-1}$; $HI=Y/B=1/2.4=0.42$ \ans $TE=B/T$ with $T\approx0.8ET_c$: \[36.0\times1000/(0.8\times506\times10)=8.89\ \text{kg\,m$^{-3}$}\quad\text{\finans}\]

}{}
\IfFileExists{extra/teny10.tex}{\subsection{Worked solutions on the ten soils (third scenario)}
The same questions are repeated for different working conditions, to show how the answer moves with them.

\exbox{Ten-soil example 10c: Coarse loam (S01)}{Coarse loam: genotype G4 with $Y_m=11.5$ Mg\,ha$^{-1}$, $K_y=1.28$ is grown in the 125-day season on $TAW=90$ mm. Find the relative yield, the yield, the water productivity and the biomass if the harvest index is 0.43.}\sol $ET_a/ET_c=711/749=0.950$; $Y/Y_m=1-1.28(1-0.950)=0.936$ \ans $Y=10.76$ Mg\,ha$^{-1}$ \ans $WP=0.15$ kg\,m$^{-3}$ \ans Biomass $B=Y/HI=10.76/0.43=25.0$ Mg\,ha$^{-1}$ \ans

\exbox{Ten-soil example 10c: Coarse sand (S02)}{Coarse sand: genotype G4 with $Y_m=11.5$ Mg\,ha$^{-1}$, $K_y=1.28$ is grown in the 125-day season on $TAW=30$ mm. Find the relative yield, the yield, the water productivity and the biomass if the harvest index is 0.43.}\sol $ET_a/ET_c=716/749=0.956$; $Y/Y_m=1-1.28(1-0.956)=0.944$ \ans $Y=10.86$ Mg\,ha$^{-1}$ \ans $WP=0.15$ kg\,m$^{-3}$ \ans Biomass $B=Y/HI=10.86/0.43=25.3$ Mg\,ha$^{-1}$ \ans

\exbox{Ten-soil example 10c: Sand (S03)}{Sand: genotype G4 with $Y_m=11.5$ Mg\,ha$^{-1}$, $K_y=1.28$ is grown in the 125-day season on $TAW=36$ mm. Find the relative yield, the yield, the water productivity and the biomass if the harvest index is 0.43.}\sol $ET_a/ET_c=712/749=0.951$; $Y/Y_m=1-1.28(1-0.951)=0.937$ \ans $Y=10.78$ Mg\,ha$^{-1}$ \ans $WP=0.15$ kg\,m$^{-3}$ \ans Biomass $B=Y/HI=10.78/0.43=25.1$ Mg\,ha$^{-1}$ \ans

\exbox{Ten-soil example 10c: Loamy sand (S04)}{Loamy sand: genotype G4 with $Y_m=11.5$ Mg\,ha$^{-1}$, $K_y=1.28$ is grown in the 125-day season on $TAW=54$ mm. Find the relative yield, the yield, the water productivity and the biomass if the harvest index is 0.43.}\sol $ET_a/ET_c=714/749=0.954$; $Y/Y_m=1-1.28(1-0.954)=0.941$ \ans $Y=10.82$ Mg\,ha$^{-1}$ \ans $WP=0.15$ kg\,m$^{-3}$ \ans Biomass $B=Y/HI=10.82/0.43=25.2$ Mg\,ha$^{-1}$ \ans

\exbox{Ten-soil example 10c: Sandy loam (S05)}{Sandy loam: genotype G4 with $Y_m=11.5$ Mg\,ha$^{-1}$, $K_y=1.28$ is grown in the 125-day season on $TAW=72$ mm. Find the relative yield, the yield, the water productivity and the biomass if the harvest index is 0.43.}\sol $ET_a/ET_c=701/749=0.936$; $Y/Y_m=1-1.28(1-0.936)=0.918$ \ans $Y=10.56$ Mg\,ha$^{-1}$ \ans $WP=0.15$ kg\,m$^{-3}$ \ans Biomass $B=Y/HI=10.56/0.43=24.6$ Mg\,ha$^{-1}$ \ans

\exbox{Ten-soil example 10c: Loam (S06)}{Loam: genotype G4 with $Y_m=11.5$ Mg\,ha$^{-1}$, $K_y=1.28$ is grown in the 125-day season on $TAW=102$ mm. Find the relative yield, the yield, the water productivity and the biomass if the harvest index is 0.43.}\sol $ET_a/ET_c=701/749=0.937$; $Y/Y_m=1-1.28(1-0.937)=0.919$ \ans $Y=10.57$ Mg\,ha$^{-1}$ \ans $WP=0.15$ kg\,m$^{-3}$ \ans Biomass $B=Y/HI=10.57/0.43=24.6$ Mg\,ha$^{-1}$ \ans

\exbox{Ten-soil example 10c: Silt loam (S07)}{Silt loam: genotype G4 with $Y_m=11.5$ Mg\,ha$^{-1}$, $K_y=1.28$ is grown in the 125-day season on $TAW=120$ mm. Find the relative yield, the yield, the water productivity and the biomass if the harvest index is 0.43.}\sol $ET_a/ET_c=705/749=0.941$; $Y/Y_m=1-1.28(1-0.941)=0.925$ \ans $Y=10.64$ Mg\,ha$^{-1}$ \ans $WP=0.15$ kg\,m$^{-3}$ \ans Biomass $B=Y/HI=10.64/0.43=24.7$ Mg\,ha$^{-1}$ \ans

\exbox{Ten-soil example 10c: Sandy clay loam (S08)}{Sandy clay loam: genotype G4 with $Y_m=11.5$ Mg\,ha$^{-1}$, $K_y=1.28$ is grown in the 125-day season on $TAW=84$ mm. Find the relative yield, the yield, the water productivity and the biomass if the harvest index is 0.43.}\sol $ET_a/ET_c=705/749=0.941$; $Y/Y_m=1-1.28(1-0.941)=0.925$ \ans $Y=10.63$ Mg\,ha$^{-1}$ \ans $WP=0.15$ kg\,m$^{-3}$ \ans Biomass $B=Y/HI=10.63/0.43=24.7$ Mg\,ha$^{-1}$ \ans

\exbox{Ten-soil example 10c: Clay loam (S09)}{Clay loam: genotype G4 with $Y_m=11.5$ Mg\,ha$^{-1}$, $K_y=1.28$ is grown in the 125-day season on $TAW=96$ mm. Find the relative yield, the yield, the water productivity and the biomass if the harvest index is 0.43.}\sol $ET_a/ET_c=705/749=0.942$; $Y/Y_m=1-1.28(1-0.942)=0.925$ \ans $Y=10.64$ Mg\,ha$^{-1}$ \ans $WP=0.15$ kg\,m$^{-3}$ \ans Biomass $B=Y/HI=10.64/0.43=24.7$ Mg\,ha$^{-1}$ \ans

\exbox{Ten-soil example 10c: Silty clay loam (S10)}{Silty clay loam: genotype G4 with $Y_m=11.5$ Mg\,ha$^{-1}$, $K_y=1.28$ is grown in the 125-day season on $TAW=102$ mm. Find the relative yield, the yield, the water productivity and the biomass if the harvest index is 0.43.}\sol $ET_a/ET_c=701/749=0.937$; $Y/Y_m=1-1.28(1-0.937)=0.919$ \ans $Y=10.57$ Mg\,ha$^{-1}$ \ans $WP=0.15$ kg\,m$^{-3}$ \ans Biomass $B=Y/HI=10.57/0.43=24.6$ Mg\,ha$^{-1}$ \ans

}{}
\IfFileExists{extra/tenby10.tex}{\subsection{Worked solutions on ten crops (hot climate)}
The ten crops of the library differ in their demand, their roots and their sensitivity. The same type of question is answered for each of them at $ET_0=7.5$ mm\,d$^{-1}$; the data are typical values of the crops (Doorenbos--Kassam and FAO 56).

\exbox{Crop example 10c: Maize}{Maize: with $Y_m=11.5$ Mg\,ha$^{-1}$ and a seasonal $ET_c\approx900$ mm find the potential water productivity and the harvest index implied by a biomass of 2.4 times the yield.}\sol \givenline{$Y_m=11.5$ Mg\,ha$^{-1}$}\ans $WP=Y_m/ET_c=11.5\times1000/(900\times10)=1.28$ kg\,m$^{-3}$ \ans $B=2.4\times11.5=27.6$ Mg\,ha$^{-1}$; $HI=Y/B=1/2.4=0.42$ \ans $TE=B/T$ with $T\approx0.8ET_c$: $27.6\times1000/(0.8\times900\times10)=3.83$ kg\,m$^{-3}$ \finans

\exbox{Crop example 10c: Wheat}{Wheat: with $Y_m=8.0$ Mg\,ha$^{-1}$ and a seasonal $ET_c\approx1035$ mm find the potential water productivity and the harvest index implied by a biomass of 2.4 times the yield.}\sol \givenline{$Y_m=8.0$ Mg\,ha$^{-1}$}\ans $WP=Y_m/ET_c=8.0\times1000/(1035\times10)=0.77$ kg\,m$^{-3}$ \ans $B=2.4\times8.0=19.2$ Mg\,ha$^{-1}$; $HI=Y/B=1/2.4=0.42$ \ans $TE=B/T$ with $T\approx0.8ET_c$: $19.2\times1000/(0.8\times1035\times10)=2.32$ kg\,m$^{-3}$ \finans

\exbox{Crop example 10c: Cotton}{Cotton: with $Y_m=4.5$ Mg\,ha$^{-1}$ and a seasonal $ET_c\approx1242$ mm find the potential water productivity and the harvest index implied by a biomass of 2.4 times the yield.}\sol \givenline{$Y_m=4.5$ Mg\,ha$^{-1}$}\ans $WP=Y_m/ET_c=4.5\times1000/(1242\times10)=0.36$ kg\,m$^{-3}$ \ans $B=2.4\times4.5=10.8$ Mg\,ha$^{-1}$; $HI=Y/B=1/2.4=0.42$ \ans $TE=B/T$ with $T\approx0.8ET_c$: $10.8\times1000/(0.8\times1242\times10)=1.09$ kg\,m$^{-3}$ \finans

\exbox{Crop example 10c: Tomato}{Tomato: with $Y_m=75.0$ Mg\,ha$^{-1}$ and a seasonal $ET_c\approx932$ mm find the potential water productivity and the harvest index implied by a biomass of 2.4 times the yield.}\sol \givenline{$Y_m=75.0$ Mg\,ha$^{-1}$}\ans $WP=Y_m/ET_c=75.0\times1000/(932\times10)=8.05$ kg\,m$^{-3}$ \ans $B=2.4\times75.0=180.0$ Mg\,ha$^{-1}$; $HI=Y/B=1/2.4=0.42$ \ans $TE=B/T$ with $T\approx0.8ET_c$: $180.0\times1000/(0.8\times932\times10)=24.15$ kg\,m$^{-3}$ \finans

\exbox{Crop example 10c: Potato}{Potato: with $Y_m=40.0$ Mg\,ha$^{-1}$ and a seasonal $ET_c\approx897$ mm find the potential water productivity and the harvest index implied by a biomass of 2.4 times the yield.}\sol \givenline{$Y_m=40.0$ Mg\,ha$^{-1}$}\ans $WP=Y_m/ET_c=40.0\times1000/(897\times10)=4.46$ kg\,m$^{-3}$ \ans $B=2.4\times40.0=96.0$ Mg\,ha$^{-1}$; $HI=Y/B=1/2.4=0.42$ \ans $TE=B/T$ with $T\approx0.8ET_c$: $96.0\times1000/(0.8\times897\times10)=13.38$ kg\,m$^{-3}$ \finans

\exbox{Crop example 10c: Soybean}{Soybean: with $Y_m=3.5$ Mg\,ha$^{-1}$ and a seasonal $ET_c\approx897$ mm find the potential water productivity and the harvest index implied by a biomass of 2.4 times the yield.}\sol \givenline{$Y_m=3.5$ Mg\,ha$^{-1}$}\ans $WP=Y_m/ET_c=3.5\times1000/(897\times10)=0.39$ kg\,m$^{-3}$ \ans $B=2.4\times3.5=8.4$ Mg\,ha$^{-1}$; $HI=Y/B=1/2.4=0.42$ \ans $TE=B/T$ with $T\approx0.8ET_c$: $8.4\times1000/(0.8\times897\times10)=1.17$ kg\,m$^{-3}$ \finans

\exbox{Crop example 10c: Sugar beet}{Sugar beet: with $Y_m=60.0$ Mg\,ha$^{-1}$ and a seasonal $ET_c\approx1296$ mm find the potential water productivity and the harvest index implied by a biomass of 2.4 times the yield.}\sol \givenline{$Y_m=60.0$ Mg\,ha$^{-1}$}\ans $WP=Y_m/ET_c=60.0\times1000/(1296\times10)=4.63$ kg\,m$^{-3}$ \ans $B=2.4\times60.0=144.0$ Mg\,ha$^{-1}$; $HI=Y/B=1/2.4=0.42$ \ans $TE=B/T$ with $T\approx0.8ET_c$: $144.0\times1000/(0.8\times1296\times10)=13.89$ kg\,m$^{-3}$ \finans

\exbox{Crop example 10c: Sunflower}{Sunflower: with $Y_m=3.0$ Mg\,ha$^{-1}$ and a seasonal $ET_c\approx825$ mm find the potential water productivity and the harvest index implied by a biomass of 2.4 times the yield.}\sol \givenline{$Y_m=3.0$ Mg\,ha$^{-1}$}\ans $WP=Y_m/ET_c=3.0\times1000/(825\times10)=0.36$ kg\,m$^{-3}$ \ans $B=2.4\times3.0=7.2$ Mg\,ha$^{-1}$; $HI=Y/B=1/2.4=0.42$ \ans $TE=B/T$ with $T\approx0.8ET_c$: $7.2\times1000/(0.8\times825\times10)=1.09$ kg\,m$^{-3}$ \finans

\exbox{Crop example 10c: Onion}{Onion: with $Y_m=50.0$ Mg\,ha$^{-1}$ and a seasonal $ET_c\approx945$ mm find the potential water productivity and the harvest index implied by a biomass of 2.4 times the yield.}\sol \givenline{$Y_m=50.0$ Mg\,ha$^{-1}$}\ans $WP=Y_m/ET_c=50.0\times1000/(945\times10)=5.29$ kg\,m$^{-3}$ \ans $B=2.4\times50.0=120.0$ Mg\,ha$^{-1}$; $HI=Y/B=1/2.4=0.42$ \ans $TE=B/T$ with $T\approx0.8ET_c$: $120.0\times1000/(0.8\times945\times10)=15.87$ kg\,m$^{-3}$ \finans

\exbox{Crop example 10c: Alfalfa}{Alfalfa: with $Y_m=15.0$ Mg\,ha$^{-1}$ and a seasonal $ET_c\approx690$ mm find the potential water productivity and the harvest index implied by a biomass of 2.4 times the yield.}\sol \givenline{$Y_m=15.0$ Mg\,ha$^{-1}$}\ans $WP=Y_m/ET_c=15.0\times1000/(690\times10)=2.17$ kg\,m$^{-3}$ \ans $B=2.4\times15.0=36.0$ Mg\,ha$^{-1}$; $HI=Y/B=1/2.4=0.42$ \ans $TE=B/T$ with $T\approx0.8ET_c$: $36.0\times1000/(0.8\times690\times10)=6.52$ kg\,m$^{-3}$ \finans

}{}
\IfFileExists{extra/tenx10.tex}{\subsection{Worked solutions on the ten soils (second scenario)}
The same questions are repeated for different working conditions, to show how the answer moves with them.

\exbox{Ten-soil example 10b: Coarse loam (S01)}{Coarse loam: genotype G4 with $Y_m=11.5$ Mg\,ha$^{-1}$, $K_y=1.28$ is grown in the 125-day season on $TAW=90$ mm. Find the relative yield, the yield, the water productivity and the biomass if the harvest index is 0.43.}\sol \givenline{$Y_m=11.5$ Mg\,ha$^{-1}$ ; $K_y=1.28$ ; $TAW=90$ mm}\ans $ET_a/ET_c=759/759=1.000$; $Y/Y_m=1-1.28(1-1.000)=1.000$ \ans $Y=11.50$ Mg\,ha$^{-1}$ \ans $WP=0.15$ kg\,m$^{-3}$ \ans Biomass $B=Y/HI=11.50/0.43=26.7$ Mg\,ha$^{-1}$ \finans

\exbox{Ten-soil example 10b: Coarse sand (S02)}{Coarse sand: genotype G4 with $Y_m=11.5$ Mg\,ha$^{-1}$, $K_y=1.28$ is grown in the 125-day season on $TAW=30$ mm. Find the relative yield, the yield, the water productivity and the biomass if the harvest index is 0.43.}\sol \givenline{$Y_m=11.5$ Mg\,ha$^{-1}$ ; $K_y=1.28$ ; $TAW=30$ mm}\ans $ET_a/ET_c=759/759=1.000$; $Y/Y_m=1-1.28(1-1.000)=1.000$ \ans $Y=11.50$ Mg\,ha$^{-1}$ \ans $WP=0.15$ kg\,m$^{-3}$ \ans Biomass $B=Y/HI=11.50/0.43=26.7$ Mg\,ha$^{-1}$ \finans

\exbox{Ten-soil example 10b: Sand (S03)}{Sand: genotype G4 with $Y_m=11.5$ Mg\,ha$^{-1}$, $K_y=1.28$ is grown in the 125-day season on $TAW=36$ mm. Find the relative yield, the yield, the water productivity and the biomass if the harvest index is 0.43.}\sol \givenline{$Y_m=11.5$ Mg\,ha$^{-1}$ ; $K_y=1.28$ ; $TAW=36$ mm}\ans $ET_a/ET_c=759/759=1.000$; $Y/Y_m=1-1.28(1-1.000)=1.000$ \ans $Y=11.50$ Mg\,ha$^{-1}$ \ans $WP=0.15$ kg\,m$^{-3}$ \ans Biomass $B=Y/HI=11.50/0.43=26.7$ Mg\,ha$^{-1}$ \finans

\exbox{Ten-soil example 10b: Loamy sand (S04)}{Loamy sand: genotype G4 with $Y_m=11.5$ Mg\,ha$^{-1}$, $K_y=1.28$ is grown in the 125-day season on $TAW=54$ mm. Find the relative yield, the yield, the water productivity and the biomass if the harvest index is 0.43.}\sol \givenline{$Y_m=11.5$ Mg\,ha$^{-1}$ ; $K_y=1.28$ ; $TAW=54$ mm}\ans $ET_a/ET_c=759/759=1.000$; $Y/Y_m=1-1.28(1-1.000)=1.000$ \ans $Y=11.50$ Mg\,ha$^{-1}$ \ans $WP=0.15$ kg\,m$^{-3}$ \ans Biomass $B=Y/HI=11.50/0.43=26.7$ Mg\,ha$^{-1}$ \finans

\exbox{Ten-soil example 10b: Sandy loam (S05)}{Sandy loam: genotype G4 with $Y_m=11.5$ Mg\,ha$^{-1}$, $K_y=1.28$ is grown in the 125-day season on $TAW=72$ mm. Find the relative yield, the yield, the water productivity and the biomass if the harvest index is 0.43.}\sol \givenline{$Y_m=11.5$ Mg\,ha$^{-1}$ ; $K_y=1.28$ ; $TAW=72$ mm}\ans $ET_a/ET_c=759/759=1.000$; $Y/Y_m=1-1.28(1-1.000)=1.000$ \ans $Y=11.50$ Mg\,ha$^{-1}$ \ans $WP=0.15$ kg\,m$^{-3}$ \ans Biomass $B=Y/HI=11.50/0.43=26.7$ Mg\,ha$^{-1}$ \finans

\exbox{Ten-soil example 10b: Loam (S06)}{Loam: genotype G4 with $Y_m=11.5$ Mg\,ha$^{-1}$, $K_y=1.28$ is grown in the 125-day season on $TAW=102$ mm. Find the relative yield, the yield, the water productivity and the biomass if the harvest index is 0.43.}\sol \givenline{$Y_m=11.5$ Mg\,ha$^{-1}$ ; $K_y=1.28$ ; $TAW=102$ mm}\ans $ET_a/ET_c=759/759=1.000$; $Y/Y_m=1-1.28(1-1.000)=1.000$ \ans $Y=11.50$ Mg\,ha$^{-1}$ \ans $WP=0.15$ kg\,m$^{-3}$ \ans Biomass $B=Y/HI=11.50/0.43=26.7$ Mg\,ha$^{-1}$ \finans

\exbox{Ten-soil example 10b: Silt loam (S07)}{Silt loam: genotype G4 with $Y_m=11.5$ Mg\,ha$^{-1}$, $K_y=1.28$ is grown in the 125-day season on $TAW=120$ mm. Find the relative yield, the yield, the water productivity and the biomass if the harvest index is 0.43.}\sol \givenline{$Y_m=11.5$ Mg\,ha$^{-1}$ ; $K_y=1.28$ ; $TAW=120$ mm}\ans $ET_a/ET_c=759/759=1.000$; $Y/Y_m=1-1.28(1-1.000)=1.000$ \ans $Y=11.50$ Mg\,ha$^{-1}$ \ans $WP=0.15$ kg\,m$^{-3}$ \ans Biomass $B=Y/HI=11.50/0.43=26.7$ Mg\,ha$^{-1}$ \finans

\exbox{Ten-soil example 10b: Sandy clay loam (S08)}{Sandy clay loam: genotype G4 with $Y_m=11.5$ Mg\,ha$^{-1}$, $K_y=1.28$ is grown in the 125-day season on $TAW=84$ mm. Find the relative yield, the yield, the water productivity and the biomass if the harvest index is 0.43.}\sol \givenline{$Y_m=11.5$ Mg\,ha$^{-1}$ ; $K_y=1.28$ ; $TAW=84$ mm}\ans $ET_a/ET_c=759/759=1.000$; $Y/Y_m=1-1.28(1-1.000)=1.000$ \ans $Y=11.50$ Mg\,ha$^{-1}$ \ans $WP=0.15$ kg\,m$^{-3}$ \ans Biomass $B=Y/HI=11.50/0.43=26.7$ Mg\,ha$^{-1}$ \finans

\exbox{Ten-soil example 10b: Clay loam (S09)}{Clay loam: genotype G4 with $Y_m=11.5$ Mg\,ha$^{-1}$, $K_y=1.28$ is grown in the 125-day season on $TAW=96$ mm. Find the relative yield, the yield, the water productivity and the biomass if the harvest index is 0.43.}\sol \givenline{$Y_m=11.5$ Mg\,ha$^{-1}$ ; $K_y=1.28$ ; $TAW=96$ mm}\ans $ET_a/ET_c=759/759=1.000$; $Y/Y_m=1-1.28(1-1.000)=1.000$ \ans $Y=11.50$ Mg\,ha$^{-1}$ \ans $WP=0.15$ kg\,m$^{-3}$ \ans Biomass $B=Y/HI=11.50/0.43=26.7$ Mg\,ha$^{-1}$ \finans

\exbox{Ten-soil example 10b: Silty clay loam (S10)}{Silty clay loam: genotype G4 with $Y_m=11.5$ Mg\,ha$^{-1}$, $K_y=1.28$ is grown in the 125-day season on $TAW=102$ mm. Find the relative yield, the yield, the water productivity and the biomass if the harvest index is 0.43.}\sol \givenline{$Y_m=11.5$ Mg\,ha$^{-1}$ ; $K_y=1.28$ ; $TAW=102$ mm}\ans $ET_a/ET_c=759/759=1.000$; $Y/Y_m=1-1.28(1-1.000)=1.000$ \ans $Y=11.50$ Mg\,ha$^{-1}$ \ans $WP=0.15$ kg\,m$^{-3}$ \ans Biomass $B=Y/HI=11.50/0.43=26.7$ Mg\,ha$^{-1}$ \finans

}{}
\IfFileExists{extra/tenbx10.tex}{\subsection{Worked solutions on ten crops (cool climate)}
The ten crops of the library differ in their demand, their roots and their sensitivity. The same type of question is answered for each of them at $ET_0=4.0$ mm\,d$^{-1}$; the data are typical values of the crops (Doorenbos--Kassam and FAO 56).

\exbox{Crop example 10b: Maize}{Maize: with $Y_m=11.5$ Mg\,ha$^{-1}$ and a seasonal $ET_c\approx480$ mm find the potential water productivity and the harvest index implied by a biomass of 2.4 times the yield.}\sol \givenline{$Y_m=11.5$ Mg\,ha$^{-1}$}\ans $WP=Y_m/ET_c=11.5\times1000/(480\times10)=2.40$ kg\,m$^{-3}$ \ans $B=2.4\times11.5=27.6$ Mg\,ha$^{-1}$; $HI=Y/B=1/2.4=0.42$ \ans $TE=B/T$ with $T\approx0.8ET_c$: $27.6\times1000/(0.8\times480\times10)=7.19$ kg\,m$^{-3}$ \finans

\exbox{Crop example 10b: Wheat}{Wheat: with $Y_m=8.0$ Mg\,ha$^{-1}$ and a seasonal $ET_c\approx552$ mm find the potential water productivity and the harvest index implied by a biomass of 2.4 times the yield.}\sol \givenline{$Y_m=8.0$ Mg\,ha$^{-1}$}\ans $WP=Y_m/ET_c=8.0\times1000/(552\times10)=1.45$ kg\,m$^{-3}$ \ans $B=2.4\times8.0=19.2$ Mg\,ha$^{-1}$; $HI=Y/B=1/2.4=0.42$ \ans $TE=B/T$ with $T\approx0.8ET_c$: $19.2\times1000/(0.8\times552\times10)=4.35$ kg\,m$^{-3}$ \finans

\exbox{Crop example 10b: Cotton}{Cotton: with $Y_m=4.5$ Mg\,ha$^{-1}$ and a seasonal $ET_c\approx662$ mm find the potential water productivity and the harvest index implied by a biomass of 2.4 times the yield.}\sol \givenline{$Y_m=4.5$ Mg\,ha$^{-1}$}\ans $WP=Y_m/ET_c=4.5\times1000/(662\times10)=0.68$ kg\,m$^{-3}$ \ans $B=2.4\times4.5=10.8$ Mg\,ha$^{-1}$; $HI=Y/B=1/2.4=0.42$ \ans $TE=B/T$ with $T\approx0.8ET_c$: $10.8\times1000/(0.8\times662\times10)=2.04$ kg\,m$^{-3}$ \finans

\exbox{Crop example 10b: Tomato}{Tomato: with $Y_m=75.0$ Mg\,ha$^{-1}$ and a seasonal $ET_c\approx497$ mm find the potential water productivity and the harvest index implied by a biomass of 2.4 times the yield.}\sol \givenline{$Y_m=75.0$ Mg\,ha$^{-1}$}\ans $WP=Y_m/ET_c=75.0\times1000/(497\times10)=15.10$ kg\,m$^{-3}$ \ans $B=2.4\times75.0=180.0$ Mg\,ha$^{-1}$; $HI=Y/B=1/2.4=0.42$ \ans $TE=B/T$ with $T\approx0.8ET_c$: $180.0\times1000/(0.8\times497\times10)=45.29$ kg\,m$^{-3}$ \finans

\exbox{Crop example 10b: Potato}{Potato: with $Y_m=40.0$ Mg\,ha$^{-1}$ and a seasonal $ET_c\approx478$ mm find the potential water productivity and the harvest index implied by a biomass of 2.4 times the yield.}\sol \givenline{$Y_m=40.0$ Mg\,ha$^{-1}$}\ans $WP=Y_m/ET_c=40.0\times1000/(478\times10)=8.36$ kg\,m$^{-3}$ \ans $B=2.4\times40.0=96.0$ Mg\,ha$^{-1}$; $HI=Y/B=1/2.4=0.42$ \ans $TE=B/T$ with $T\approx0.8ET_c$: $96.0\times1000/(0.8\times478\times10)=25.08$ kg\,m$^{-3}$ \finans

\exbox{Crop example 10b: Soybean}{Soybean: with $Y_m=3.5$ Mg\,ha$^{-1}$ and a seasonal $ET_c\approx478$ mm find the potential water productivity and the harvest index implied by a biomass of 2.4 times the yield.}\sol \givenline{$Y_m=3.5$ Mg\,ha$^{-1}$}\ans $WP=Y_m/ET_c=3.5\times1000/(478\times10)=0.73$ kg\,m$^{-3}$ \ans $B=2.4\times3.5=8.4$ Mg\,ha$^{-1}$; $HI=Y/B=1/2.4=0.42$ \ans $TE=B/T$ with $T\approx0.8ET_c$: $8.4\times1000/(0.8\times478\times10)=2.19$ kg\,m$^{-3}$ \finans

\exbox{Crop example 10b: Sugar beet}{Sugar beet: with $Y_m=60.0$ Mg\,ha$^{-1}$ and a seasonal $ET_c\approx691$ mm find the potential water productivity and the harvest index implied by a biomass of 2.4 times the yield.}\sol \givenline{$Y_m=60.0$ Mg\,ha$^{-1}$}\ans $WP=Y_m/ET_c=60.0\times1000/(691\times10)=8.68$ kg\,m$^{-3}$ \ans $B=2.4\times60.0=144.0$ Mg\,ha$^{-1}$; $HI=Y/B=1/2.4=0.42$ \ans $TE=B/T$ with $T\approx0.8ET_c$: $144.0\times1000/(0.8\times691\times10)=26.04$ kg\,m$^{-3}$ \finans

\exbox{Crop example 10b: Sunflower}{Sunflower: with $Y_m=3.0$ Mg\,ha$^{-1}$ and a seasonal $ET_c\approx440$ mm find the potential water productivity and the harvest index implied by a biomass of 2.4 times the yield.}\sol \givenline{$Y_m=3.0$ Mg\,ha$^{-1}$}\ans $WP=Y_m/ET_c=3.0\times1000/(440\times10)=0.68$ kg\,m$^{-3}$ \ans $B=2.4\times3.0=7.2$ Mg\,ha$^{-1}$; $HI=Y/B=1/2.4=0.42$ \ans $TE=B/T$ with $T\approx0.8ET_c$: $7.2\times1000/(0.8\times440\times10)=2.05$ kg\,m$^{-3}$ \finans

\exbox{Crop example 10b: Onion}{Onion: with $Y_m=50.0$ Mg\,ha$^{-1}$ and a seasonal $ET_c\approx504$ mm find the potential water productivity and the harvest index implied by a biomass of 2.4 times the yield.}\sol \givenline{$Y_m=50.0$ Mg\,ha$^{-1}$}\ans $WP=Y_m/ET_c=50.0\times1000/(504\times10)=9.92$ kg\,m$^{-3}$ \ans $B=2.4\times50.0=120.0$ Mg\,ha$^{-1}$; $HI=Y/B=1/2.4=0.42$ \ans $TE=B/T$ with $T\approx0.8ET_c$: $120.0\times1000/(0.8\times504\times10)=29.76$ kg\,m$^{-3}$ \finans

\exbox{Crop example 10b: Alfalfa}{Alfalfa: with $Y_m=15.0$ Mg\,ha$^{-1}$ and a seasonal $ET_c\approx368$ mm find the potential water productivity and the harvest index implied by a biomass of 2.4 times the yield.}\sol \givenline{$Y_m=15.0$ Mg\,ha$^{-1}$}\ans $WP=Y_m/ET_c=15.0\times1000/(368\times10)=4.08$ kg\,m$^{-3}$ \ans $B=2.4\times15.0=36.0$ Mg\,ha$^{-1}$; $HI=Y/B=1/2.4=0.42$ \ans $TE=B/T$ with $T\approx0.8ET_c$: $36.0\times1000/(0.8\times368\times10)=12.23$ kg\,m$^{-3}$ \finans

}{}
\IfFileExists{extra/ex10.tex}{\backhead{EXERCISES}{khred}
\begin{exlist}
\item WP for 9.5 Mg/ha at 460 mm.\exans{2.07}
\item CWSI for $T_c=31$, $T_a=34$ with limits $-3$, $+4$.\exans{0.00}
\item HI for grain 5.2, biomass 12.\exans{0.43}
\item Error d.f. for 12 genotypes, 3 blocks.\exans{22}
\item LSD for $s=0.8$, $r=3$, $t=2.1$.\exans{1.37}
\item Biomass from $V=0.25$ m$^3$, $\rho_v=15$.\exans{3.75 kg}
\item TE for $B=1.8$ kg, $T=720$ L.\exans{2.5 g/kg}
\item Why randomise blocks?\exans{to remove gradients}
\item State two traits of a drought-tolerant genotype.\exans{low CWSI; high WP}
\item What does a CWSI of 1 mean?\exans{fully stressed}
\end{exlist}
}{}
\IfFileExists{extra/dyk10.tex}{\dykband{DO YOU KNOW ?}
\begin{exlist}
\item Why compare genotypes on the same water?
\item What is the harvest index?
\item How is canopy temperature turned into a stress index?
\item Why is a randomised complete block used?
\item What does the volumetric scan measure?
\item What limits the use of WP alone?
\item What is transpiration efficiency?
\item How many plots has 20 genotypes by 4 blocks?
\end{exlist}
}{}
\IfFileExists{extra/obj10.tex}{\dykband{OBJECTIVE TYPE QUESTIONS}
\begin{exlist}
\item WP is\opt{$Y/ET_a$}{$ET_a/Y$}{$Y\cdot ET$}{$B/Y$}
\item A CWSI near 0\opt{no stress}{full stress}{dead crop}{no sun}
\item HI is\opt{grain/biomass}{biomass/grain}{water/grain}{none}
\item Blocks remove\opt{gradients}{genotypes}{weather}{yield}
\item Hot canopy means\opt{stress}{no stress}{rain}{nothing}
\item LSD uses\opt{$t$ and $s$}{wind}{slope}{RH}
\item Scan returns\opt{volume}{price}{rain}{salt}
\item TE is\opt{$B/T$}{$T/B$}{$Y/I$}{$P/E$}
\item Plots = g times\opt{r}{s}{t}{p}
\item Best for water-scarce schemes\opt{highest WP}{highest ET}{lowest HI}{largest canopy}
\end{exlist}
\objkey{1.~(a)\quad 2.~(a)\quad 3.~(a)\quad 4.~(a)\quad 5.~(a)\quad 6.~(a)\quad 7.~(a)\quad 8.~(a)\quad 9.~(a)\quad 10.~(a)}
}{}
